% 本文件由 LaTeX 排版 Agent 根据有效 translation.json 自动生成。
% author=Gregory of Nyssa; work_id=2914; title=论人的受造

\chapter{论人的受造}

% structure: p0001 source=h1 target=omitted reason=page-title-already-rendered-by-parent

% structure: p0002 source=h2 target=section reason=relative-heading-rank; base=h2

\section{注释}

这部著作旨在补充并完成\Person{圣巴西略}的\WorkTitle{六天创世}，并且预设读者熟悉该论著。创造世界的叙述未被详尽讨论：它被提及，但主要是为了强调这样的观念：世界被预备为人类统治的领域。另一方面，\Person{额我略}表明，人是\Quote{慎重地}被造的，在本性上适合统治其他受造物，按\Quote{天主的肖像}被造，拥有各种道德属性，并拥有理性，同时与天主性区别在于人的心智藉感官接收信息，并依赖感官知觉外在事物。身体被造为心智的工具，适应于理性存在的使用；正是由于拥有\Quote{理性灵魂}，以及\Quote{自然}或\Quote{生长性}和\Quote{感性}灵魂，人才与低等动物区别开来。同时，他的心智藉感官运作：其本性无法理解（在此相似于它所肖像的天主性），并且它与身体的关系被相当详尽地讨论（第十二至十五章）。心智与身体的联系是无法言喻的：不能通过假设心智居于身体的任何特定部分来解释：心智作用于整个身体，也被整个身体作用，依赖于物质和形体本性以获取知觉元素，因此知觉既需要身体也需要心智。但正是理性元素才正确属于\Quote{灵魂}之名：营养和感性官能只是从高于自身者那里借用名称。人首先被造为\Quote{天主的肖像}；这一概念排除性的区别。在人的首次创造中，按照天主的预知，所有人类都包含在内：\Quote{我们的整个本性从起初直到末后}就是\Quote{那自有者的一个肖像}。若非堕落，人类的繁衍本会以天使族裔繁衍的方式发生，以某种我们不为人知的方式。人从起初地位的堕降使得藉生育的延续成为必要；正是因为这堕落及其后果呈现在天主的预知中，天主才\BibleQuote{造了他们，男的和女的}。在这方面，以及在对食物滋养的需要方面，人并非\Quote{天主的肖像}，而是显示出与低等受造物的亲缘。但这些需要并非永久：它们将随着人恢复其原先的卓越而结束（第十六至十八章）。至此，\Person{额我略}开始论述（第十九至二十章）乐园中人的食物，以及\Quote{知善恶树}。这样，在提及人的堕落之后，他转而论述人的恢复。在他看来，这源于恶的有限性；它被推迟，直到人类的总数满全。至于事物现状将以何种方式终结，我们一无所知；但它的终结是从物质的非永恒性推断出来的（第二十一至二十四章）。复活教义得到以下支持：我们知道圣经中其他事件被预言的准确性，我们在特殊情况下所经历的相似事件的经验，即那些被主复活的人，尤其是祂自己的复活。针对这种恢复不可能的论点，通过诉诸天主能力的无限性，以及从自然中观察到的相似事例的推论来回应（第二十五至二十七章）。\Person{额我略}接着处理灵魂先存的问题，拒绝这一观点，并坚持身体和灵魂一同进入存在，\Emph{潜能上}在于天主的旨意，\Emph{现实上}在于每个人藉生育而成为存在的时刻（第二十八至二十九章）。在对这最后一点的论证中，他转而详细讨论——在最后一章——人体结构；但他再次回到其主要立场：人\Quote{是作为有生命和有灵性的存在而被生育的}，而灵魂的能力逐渐在身体的物质基础上并藉此基础彰显；因此人藉灵魂较低属性的帮助而达至完美。但灵魂的真正完美不在这些属性——它们终将\Quote{被除去}——而在那构成人\Quote{天主的肖像}的较高属性。\par

% structure: p0004 source=h2 target=section reason=relative-heading-rank; base=h2

\section{绪论}

\Signature{\Person{额我略}，\Place{尼萨}主教，致其弟\Person{伯多禄}，天主的仆人。}\par

如果我们必须以金钱奖励来尊崇那些在德行上卓越的人，那么正如\Person{撒罗满}所说的，全世界金钱在平衡中与您的德行相比也显得微不足道。然而，由于欠您尊威的感恩之债远非金钱所能衡量，而神圣的复活节期要求惯常的爱之赠礼，我们便向您的高尚心灵献上，哦天主的仆人，一件礼物；它确实太小不值得献给您，但并未超出我们能力所及。礼物是一篇言论，好比一件粗糙的衣服，是从我们贫乏的才智中不无劳苦地编织而成；而这篇言论的主题，虽然可能普遍被认为大胆，却似乎并非不合适。因为唯独那位真正按天主而被造、且其灵魂被塑造为造他者的肖像的人——即我们的共同父亲和导师\Person{圣巴西略}——才相称地默想了天主的创造；他通过自己的思辨，使那崇高的宇宙秩序普遍可理解，使那由天主在真智慧中所建立的世界得以被那些藉他的理解而导向此种默观的人所认识。至于我们，甚至连相称地赞美他都做不到，却仍打算给这位伟大作家的思辨补充其所缺，不是要以插入方式篡改其著作（因为绝不可让那高贵的口舌蒙受以我们言论作权威的侮辱），而是使导师的光荣不在其弟子中暗淡。\par

因为如果，在他的\WorkTitle{六日创造}中缺乏对人的思考，那些曾是他门徒的人中没有一个认真努力去弥补这个缺陷，那么嘲笑者或许就有了攻击他伟大名声的把柄，借口他不关心在听众中培养任何智识习惯。但如今我们按自己的能力尝试阐述所缺的部分，如果我们的作品中有什么不配他的教导，那肯定要归咎于我们的老师；而如果我们的论述达不到他那崇高思辨的高度，他则免于指责，避免被认为不希望他的门徒有任何技能，而我们可能要对批评者负责，因为我们内心渺小，无法容纳导师的智慧。\par

我们拟议探讨的范围并不小：它不亚于世界上任何奇观——甚至可能比我们所知的任何奇观都更伟大，因为除人的受造物外，没有其他存在物被造成相似于天主。因此，和善的读者将很容易体谅我们所说的话，即使我们的论述远逊于主题的价值。因为我相信，我们的任务是不遗漏任何关于人的问题——无论是我们相信先前发生的事，我们现在看见的事，还是以后预期会出现的结果（因为如果当人作为默观对象时，任何与主题相关的问题被省略，我们的努力无疑会被证明未能兑现承诺）；此外，我们必须根据圣经的解释和理性的推导，将那些似乎因某种必然顺序而相互对立的关于人的陈述调和起来，以使我们的整个主题在思路和顺序上保持一致，当那些看似矛盾的陈述（如果天主的能力为无望者发现希望，为无路者发现道路）被导向同一个目的时；为了清楚起见，我认为最好按章节向你们陈述这篇论述，以便你们能简要了解整个作品中几个论证的要旨。\par

% structure: p0009 source=h2 target=section reason=relative-heading-rank; base=h2

\section{\Emph{一、其中对世界本性的部分探究，以及对先于人类受生之事物的更详细阐述}}

1. \Quote{这是天地受生的记录}，圣经说，当所有可见之物完成，每一事物各归其所，天穹环绕一切，沉重且向下倾向的物体——大地和水——相互含摄，占据了宇宙的中央；同时，作为一种维系和稳定受造物的纽带，天主的能力和智慧被植入万物的生长中，以双重作为的缰绳引导一切（因为它通过静止与运动筹划了未有之物的产生和现有之物的延续），围绕着由不运动的受造物所贡献的沉重而不变的元素，如同围绕某个固定轨道，驱策球体极其快速的运动，有如轮子，并通过它们的相互作用保持两者的不可分解性：旋转物质通过快速运动压缩周围大地的坚实身体，而那坚固且不屈的，因其不变的固定，不断增强周围旋转事物的回旋运动，并在这些运作不同的本性——我说的是静止本性与移动旋转——中产生同等程度的强度；因为大地从未从自己的基础移动，苍穹也从不减弱其猛烈或减缓其运动。\par

2. 此外，这些事物根据上智的安排，在其他事物之前被首先建造，作为整个机器的某种开端，伟大的梅瑟（我想）指示说，天主\Quote{在起初 \BibleRef{创 1:1}}创造了天地，在受造物中所有可见之物都是静止与运动的后裔，由天主的旨意产生。如今，天与地在运作上截然对立，介于对立面之间的受造物，部分分享了邻近者的性质，自身作为极端之间的中介，以致通过中介，对立面之间明显有相互接触；因为空气在某种程度上模仿火热物质的永恒运动与精微，既在其本性的轻盈性上，又在其适合运动上；然而它并非如此以至于与固体物质相疏离，因为它既不是处于不断的流动与分散中，也不是处于永久的静止状态，而是因其与每一方的亲和性，成为运作对立之间的某种边界，同时统一又分割那些本性上截然不同的事物。\par

3. 同样，液体物质也通过双重属性附属于每一对立面；就其沉重且向下倾向而言，它与土性紧密相似；但就其分有某种流动而易动的能量而言，它与运动的本性并非完全异己；由此也实现了对立面的某种混合与交汇，重量转移为运动，运动在重量中不受阻碍，以致本性上最为对立的事物相互结合，并通过作为它们之间中介的事物而彼此连接。\par

4. 但严格来说，更确切地说，相反者本身的性质并非完全没有彼此属性的掺杂，以至于，在我看来，我们在世上所见到的一切都相互协调，而创造物虽然被发现具有相反性质的各种属性，却仍与自身合一。因为运动不仅被视为位置的移动，也在变化和改变中被默观；另一方面，不动者不接受通过变化而来的运动，天主的智慧便转换了这些属性，在永恒运动者中产生了不变，而在不动者中产生了变化；这样做的目的，或许是出于上智的安排，使构成其不变性和不动性的那种自然属性，当在任何受造物中被默观时，不致让受造物被视为天主；因为可能运动或改变的事物将不再接受神性的概念。因此，大地是稳定的，却并非不变的；反之，苍天虽无变化，却也没有稳定性，为使天主的能力，通过将变化交织在稳定的性质中，并将运动与那不变者相结合，借着属性的互换，既能将两者紧密联接，又能使它们远离神性的概念；因为如前所述，这两者（无论是无常者还是可变者）都不能被认为属于更神圣的性质。\par

5. 如今，万物都已抵达各自的终点：\Quote{天和地\BibleRef{创 2:1}，}正如梅瑟所说，\Quote{都完成了，}以及介于其间的一切事物；各个事物都以其相称的美妙妆点着：苍天以星辰的光芒，海洋和空气以水栖和飞行的生物，大地以各类植物和动物，这一切都因天主的旨意而获得能力，大地同时孕育了它们；大地充满了出产，花开同时结果；草地长满了各种植物，所有山脊、峰顶、每个山坡、斜坡和洼地，都披上了嫩草，并树木的多样出产，它们刚破土而出，却立即长至完美的美丽；所有按天主的命令进入生命的野兽，我们可以想象，都在欢跃蹦跳，成群结队地按着种类在林间穿梭；每个隐蔽阴凉之处都回响着鸣禽的歌声。在海中，我们可以想象，所见景象也是如此：海洋刚因深渊汇聚而安静平息，海湾和港口天然地凹入海岸，使海洋与陆地和解；微波轻漾，在美丽上与草地争妍，在轻拂海面的和风下泛起涟漪；天地间的所有受造财富都已备妥，却无人分享。\par

% structure: p0015 source=h2 target=section reason=relative-heading-rank; base=h2

\section{II. \Emph{为何人在创造之后最后出现}}

1. 因为那伟大而宝贵之物——人——尚未进入存在世界；不能指望统治者出现在其臣民之前；但当他的统治准备就绪，下一步便是君王显现。因此，万有的创造者预先为未来的君王准备了一座王室的居所（这就是大地、岛屿、海洋以及如穹顶般拱起于其上的苍天），并将各类财富储存在这座宫殿中（我所说的财富是指整个受造界，即植物和树木中的一切，以及所有具有感觉、气息和生命的事物；如果我们也把材料算作财富，那么那些因美丽而被世人视为珍贵的东西，如金银、以及你们所喜爱的宝石的物质——我说，所有这些也都隐藏在大地的怀抱中，如同在王室的宝库中），于是祂将人彰显于世间，使他成为其中一些奇观的观赏者，另一些奇观的主人；使他通过享受而认识施予者，通过所见之物的美丽和威严而追踪那超越言语和语言的造物主的能力。\par

2. 为此缘故，人在创造之后最后被引入世界，并非因其无价值而被推到末尾，而是作为一个一出生就应当成为其臣民之君王的人。正如一个好主人不会在宴席准备完毕之前将客人带进家中，而是当他做好一切适宜的准备，用适当的装饰布置好房屋、卧榻、餐桌之后，才将客人带回家，那时一切适宜款待之物已准备就绪——同样，我们本性的富足而慷慨的款待者在用各类美物装饰了居所，并准备了这盛大而多样的宴席之后，便引入了人，分派给他的任务不是获取不在场之物，而是享受在场之物；为此，祂将双重构造的本能作为基础赐予他，将神性与尘土相混合，以便借着两者，他能自然而适当地趋向每一种享受：借着他更神圣的本性享受天主，借着与尘世美好事物相仿的感觉来享受它们。\par

% structure: p0018 source=h2 target=section reason=relative-heading-rank; base=h2

\section{III. \Emph{人的本性比所有可见受造界更宝贵}}

1. 但是，我们不应不经考虑就放过这一点：即世界，尽管伟大，以及其各部分，都被安置为宇宙形成的元素基础，然而宇宙的创造可以说是由天主的能力随手而成，在祂的命令下即刻存在；而人的制造则有深思熟虑作为先导。那将要存在之物由造物主在言语描述中预先显示：它应当是什么样，应当与什么样的原型相似，它为何被造，被造后它的作用是什么，它应统治什么——这些话预先考量了这一切，使人在其生成之前就已有指定的地位，并在其存在之前就拥有对万物的统治权；因为经上说：\BibleQuote{天主说：让我们照我们的肖像，按我们的模样造人，让他们管理海中的鱼、地上的走兽、天空的飞鸟、牲畜和全地。}\par

2. 何等奇妙！太阳被造，没有经过商议；同样，苍天也是如此；在这两者面前，没有任何受造物可与之相比。如此伟大的奇观仅凭一句话就形成了，而这句话并未指明何时、如何或任何此类细节。同样，在一切个别情况中：以太、星辰、中间的大气、海洋、大地、动物、植物——全凭一句话就得以存在。唯独在造人时，万有的造物主以谨慎的态度接近，以便预先为其形成预备材料，使其形态与一个原型之美相似，为他设定一个他之所以存在的目标，并为他造出一种与他的运作相适宜、相联合的本性，且适于眼前的目的。\par

% structure: p0021 source=h2 target=section reason=relative-heading-rank; base=h2

\section{IV. \Emph{人的构造处处表明其统治权}}

1. 正如在我们自己的生活中，工匠按照其用途来制作工具，同样，这位最优秀的工匠使我们的本性成为一种适合行使王权的形态，既通过灵魂的优越优势，也通过身体的形式本身，使之适于王权。因为灵魂立即显示出其王权和崇高的特征，远离卑贱的私人地位，它不为任何主人所拥有，是自治的，由其自身意志独裁地支配。因为这除了属于君王，还能属于谁呢？此外，除了这些事实之外，它是统治万有的那个本性的肖像，这无非意味着：我们的本性从一开始就被造为王权。因为，正如在人的通常用法中，制作王子肖像的人不仅塑造其形象，还通过紫色袍服来代表其王权，甚至这肖像通常被称为\Quote{国王}；同样，人的本性，既然被造来统治其余一切，也因其与万有君王的相似，被造为，可以说，一个活的肖像，与原型共享等级和名称，不穿着紫袍，也不以权杖和王冠表明其等级（因为原型本身并不穿戴这些），而是代替紫袍，披上美德——这确实是所有衣袍中最王权的；代替权杖，倚靠不朽的真福；代替王冠，佩戴正义的冠冕。这样，它在一切属于王权尊严的事上，就显示出与原型之美完全相似。\par

% structure: p0023 source=h2 target=section reason=relative-heading-rank; base=h2

\section{V. \Emph{人是天主主权的肖像}}

1. 的确，天主的美丽并不以任何形状或形体的天赋、任何颜色的美丽来装饰，而是在不可言喻的真福中被视为卓越。正如画家借助某些颜色将人的形态转移到他们的画作上，在他们的复制品上涂上相应的色调，以便原型的美丽能准确地转移到肖像上；我愿你也这样理解：我们的造物主也在描绘肖像以相似祂自己的美丽时，通过添加美德，如同用颜色一样，在我们身上显示出祂自己的主权。而且，塑造祂真正形态的色调，可以说是多样而多变的：不是红色、白色或它们的混合（无论叫什么名称），也不是描画眉毛和眼睛的一抹黑色，或通过某种组合来阴影化轮廓的凹陷，以及画家双手所构思的所有这类技艺；而是代替这些：纯洁、无欲、真福、远离一切邪恶，以及所有这类有助于在人身上形成天主相似的特质。我们的造物主正是用这样的色彩给祂自己的肖像——我们的本性——作了标记。\par

2. 如果你也考察其他表现天主美丽之点，你会发现肖像中的相似之处也完美地保存在我们身上。天主是理智和圣言：因为\BibleQuote{在起初已有圣言}，而保禄的门徒\BibleQuote{拥有基督的心意}，基督在他们内\BibleQuote{说话}。人性离此也不远：你在自己身上看到言语和悟性，是那理智和圣言的模仿。再者，天主是爱，是爱的源泉：伟大的若望如此宣告说，\BibleQuote{爱是由天主而来的}，\BibleQuote{天主是爱}\BibleRef{若一 4:7}；我们本性的塑造者也使这成为我们的特征：因为祂说：\BibleQuote{如果你们彼此相爱，众人因此就可认出你们是我的门徒}——因此，如果这爱缺失，整个肖像的印记就被改变了。天主看见并听见一切，洞察一切；你也有通过视觉和听觉认识事物的能力，以及探究和洞察事物的悟性。\par

% structure: p0026 source=h2 target=section reason=relative-heading-rank; base=h2

\section{VI. \Emph{考察理智与本性之相似：其中，以插入方式驳斥\OriginalTerm{异类者派}{Anomœans}的学说}}

1. 但愿没有人以为我是在说：天主与现有事物的接触，类似于人的操作方式，借着不同的官能。因为要在天主的单纯性中设想认知作用的多样性和差异性是不可能的；即使在我们自己身上，认知事物的官能也并不多，虽然我们通过感官以多种方式接触那些影响我们生命的事物；因为有一种官能，即内在的理智本身，它穿过每个感官器官，把握超越感官的事物：正是它，通过眼睛看见所见之物；正是它，通过听觉理解所说之事；它满足于我们所尝到的，并避开不愉快的；它用手按自己的意愿行事，借此抓取或拒绝，为此目的恰如其分地运用器官的帮助。\par

2. 那么，即使在人类中，尽管为感知目的而由自然形成的器官可能不同，但透过所有器官运作和活动，并各自适当地为了面前的对象使用它们的那一位，却是同一不变的，不因运作的差异而改变其本性；那么，人们怎会因天主的各种能力而怀疑在祂内本质的多重性呢？因为正如先知所说，\Quote{那造眼睛的}和\Quote{那栽种耳朵的}，将这些运作印刻在人性上，仿佛是相对于祂自身中的原型的显著标记，因为祂说：\Quote{让我们按我们的肖像造人 \BibleRef{创 1:26}。}\par

3. 但是，我要问，异类者的异端变成了什么？对于这段话他们会说什么？面对所引用的言辞，他们将如何捍卫其教义的空虚？他们会说，一个肖像有可能被塑造成不同的形式吗？如果圣子在本质上与圣父不相似，那么祂由不同本性所形成的相似又怎能是一个呢？因为那位说\Quote{让我们按我们的肖像造}的，并借着复数含义启示了天主圣三的那一位，如果原型彼此不相似，就不会以单数提到肖像：因为不可能将不一致的事物表现为一个相似；如果本性不同，祂肯定也会开始以不同方式塑造它们的肖像，为每个造出合适的肖像；但既然肖像是一个，而原型不是一个，那么谁的理解力如此低下，竟不知道那些与同一事物相似的事物，彼此必定相似呢？因此祂说（这话，或许，在人类生命形成之初就阻止了这种邪恶），\Quote{让我们按我们的肖像，按我们的相似造人。}\par

% structure: p0030 source=h2 target=section reason=relative-heading-rank; base=h2

\section{VII. \Emph{人为何缺乏自然的武器和覆盖物}}

1. 但他身躯的直立意味着什么？为什么那些有助于生命的力并不自然地属于他的身体？人却被带入生命，赤裸无遮盖，一无所有且贫乏，缺乏一切有用之物，按外表看来，值得怜悯而非赞叹，没有长着突出的角或尖锐的爪，也没有蹄或牙，更没有因刺而拥有的致命毒液——这些是大多数动物为防御侵害者而拥有的能力：他的身体没有毛发覆盖；然而，或许可以预期，这位被提升以统治其余受造物的，应该被自然以自身的武器所保护，以便为自己的安全无需他人的援助。然而，现在狮子、野猪、老虎、豹子及所有此类都有足够的安全自然能力：公牛有角，野兔有速度，鹿有跳跃与目光的准确，另一种野兽有庞大的身躯，其他的有长鼻，鸟有翅膀，蜜蜂有刺，一般说来，所有动物都自然铭刻着某种保护力；但唯独人，在一切之中，比那些步伐迅速的野兽更慢，比那些身躯庞大的更小，比那些被自然武器保护的更无防御；那么，有人会说，这样的存在是如何获得对万物的主权的呢？\par

2. 嗯，我认为一点也不难表明：那看似是我们的本性之缺陷的，正是我们获得对被统治受造物统治权的手段。因为如果人拥有如此的能力，能够在速度上超过马，并拥有由于坚实而不会磨损的脚，却因某种蹄或爪而增强，并携带角、刺和爪，那么这样的东西若与他的身体一同生长，他首先会成为一只野性而可怕的存在；而且，如果他不需要其臣民的合作，他就会忽略对其他受造物的统治；而现在，我们生活中必需的服务分配在我们统治下的各个动物之间，正是出于这个原因——使我们对它们的统治成为必要。\par

3. 是我们身体缓慢而困难的行动，促使马匹来供应我们的需要，并将它驯服；是我们身体的赤露，使得我们对羊群的管理成为必要，羊群以其每年的羊毛出产弥补了本性的不足；是我们从他人那里获取生活所需的事实，使役兽服从于这样的服务；此外，是我们不能像牛那样吃草的实情，使牛以其自身的劳动来服务于我们的生活，为我们减轻生活的负担；因为我们需要用牙齿和咬力来撕咬某些动物，以制服它们，狗便以其敏捷和颚部来满足我们的需要，成为了人活生生的剑；人类还发现了铁，它比突出的角或锋利的爪更强韧、更具穿透力，不像那些动物的器官那样，总是与我们一同自然生长，而是暂时与我们结盟，其余时间则安然自存；为弥补鳄鱼鳞甲的不足，有人可将那鳞甲本身用作铠甲，适时地披在皮肤上；或者，若没有，技艺也为这个目的锻造铁，它在战时服役一段时间后，又在和平时让武士不再负重；鸟类的翅膀也服务于我们的生活，使我们在巧计的帮助下，甚至不被翅膀的速度所抛离：因为其中一些被驯化，服务于捕鸟之人，另一些则被巧计所制，以满足我们的需要；此外，技艺还使我们箭矢装上羽毛，通过弓为我们的需要提供了翅膀的速度；而我们双脚在旅行中易受伤且易磨损的实情，又使受役动物的帮助成为必要：由此，我们才把鞋穿在脚上。\par

% structure: p0034 source=h2 target=section reason=relative-heading-rank; base=h2

\section{VIII. \Emph{为何人的形体是直立的；为何双手被赐予人，出于理性；其中也包含关于灵魂差异的思辨}。}

1. 然而人的形体是直立的，向上伸向天空，仰视上苍：这些都是王权的标志，显示了他的王者尊严。因为人独与众受造物如此不同，而其余一切都弯腰朝向下方，这明确指出了那些俯伏于其权下之物与那超越它们之权力之间的尊贵差异：所有其余受造物，它们身体的前部肢体都呈脚的形式，因为下弯之物需要支撑；但在人的形成中，这些肢体被造为双手，因为直立的身躯只有一个根基，稳稳地立在双脚上，便已足够其所需。\par

2. 尤其这些辅助性的双手，适应于理性的需求：事实上，若有人说双手的服务是理性本性的一种特殊属性，他并不全错；这不仅因为他的思想转向通常明显的事实，即我们通过双手在书写中的自然运用来表达我们的推理。的确，这一事实——我们通过书写来说话，并以某种方式借助双手的辅助进行交谈，通过字母的形态保存声音——并非与理性的秉赋无关；但我所说的另有其事，即双手与理性的命令协同合作。\par

3. 然而，在讨论这一点之前，让我们思考一下我们之前略过的问题（因为受造之物的秩序几乎从我们的注意中溜走）：为何生长的植物居于首位，无理性的动物紧随其后，而在这些受造之后，才出现了人？因为这或许让我们学到的，不仅是显而易见的想法——即青草在造物主看来，是为动物之故而有用，动物则是为人而造，因此，在动物之前先造了它们的食物，在人之前先造了那将服务于人类生活之物。\par

4. 但在我看来，\Person{梅瑟}通过这些事实揭示了一种隐密的教义，并暗中传授了关于灵魂的智慧，那没有这智慧的学问虽有一些想象，却无清晰的领悟。他的论述由此教导我们：生命与灵魂的能力可按三种划分来考虑。第一种只是生长和营养的能力，提供维持被滋养身体所需的适当之物，这被称为植物性的灵魂，可见于植物；因为我们在生长的植物中，可以察觉到一种缺乏感觉的活力。除此以外，还有另一种生命形式，它虽包含上述形式，此外还具有按照感觉进行管理的能力；这可见于无理性动物的本性：因为它们不仅接受营养和生长，还具有感觉和知觉的活动。但完美的身体生命，则见于理性的（我指人的）本性，它既被滋养，又具备感觉，也分有理性，并受心智的支配。\par

5. 我们可以用一种类似这样的方式来划分主题。在存在物中，一部分是理智的，一部分是有形的。我们先不讨论理智按属性的划分，因为我们的论证不涉及这些。在有形事物中，一部分完全没有生命，一部分分享生命活力。有生命的身体中，又有一部分有感觉与生命相连，一部分没有感觉；最后，有感觉的又分为有理性的和无理性的。为此，立法者说，在无生命物质（作为有生命物形态的基础）之后，这种植物性生命被造，并在植物的生长中早已存在；然后他继续引入那些由感觉支配的受造物的起源；并且，由于按照同一顺序，在那些获得肉身生命的受造物中，有感觉的可以独立于理智本性而存在，而理性原则除非与感觉混合，否则不能具有形体——为此，人被造在动物之后，恰如自然有条不紊地走向完美。因为这个理性动物——人——混合了各种形式的灵魂；他由植物性的灵魂滋养，在生长能力之上加上了感觉能力，如果我们考虑其独特本性，感觉能力处于理智本质与更物质的本质之间，比一个粗糙却比另一个精致；随后理智本质与感觉本性中精细而开明的元素发生某种结合与混合；于是人由这三者组成：正如宗徒在 \BibleBook{厄弗所}{书} 中所教导我们的，他为他们祈祷，愿他们\Quote{身体、灵魂、精神}的全部恩宠在主来临时得以保全；他用\Quote{身体}一词指营养部分，用\Quote{灵魂}指感觉部分，用\Quote{精神}指理智部分。同样，主在福音中教导经师，应当将全心、全灵、全意爱天主的诫命置于一切诫命之前：因为在这里，在我看来，这句话也指示了同样的区分，将更属肉体的存在称为\Quote{心}，将中间的存在称为\Quote{灵魂}，将更高的、理智和心智的本性称为\Quote{意}。\par

6. 同样，宗徒承认三种状态的划分，称一种为\Quote{属血肉的}，它忙于肚腹及与之相连的享乐；另一种为\Quote{属血气的}，它在德行与恶行之间持守中间立场，超越其一，却未纯粹参与另一；再一种为\Quote{属神的}，它觉察虔诚生活的完美。因此他责备格林多人沉溺于享乐和情欲时说：\Quote{你们是属血肉的\BibleRef{格前 3:3}}，不能领受更完美的道理；而在另一处，将中间类型与完美类型作比较时说：\Quote{但属血气的人不能领受圣神的事，因为在他看是愚妄；但属神的人看透万事，却没有人能看透他\BibleRef{格前 2:14}}。所以，正如属血气的人高于属血肉的人，同样，属神的人高于属血气的人。\par

7. 因此，如果圣经告诉我们，人在一切有生命之物之后被造，那么立法者无非是在向我们宣告关于灵魂的道理，即按照事物秩序中某种必然的序列，完美者最后出现：因为在理性中已包含其他各类，而在感觉中当然也存在植物性的形式，而植物性形式又只能在联系物质中被构想；因此我们可以推测，自然仿佛借助阶梯——即生命的不同特性——从低级上升到完美的形式。\par

8. 现在，既然人是理性动物，他身体的器官就必须被造成适合理性的运用；正如你可能看见乐师按乐器的形状演奏音乐，不会用竖琴吹奏，也不会用笛子弹拨，同样，我们这些器官的安排也必须适应理性，以便被发声器官撞击时，能适当发声以用于言语。为此，手被接在身体上；因为尽管我们可以数出很多日常生活中这些精巧构造的有益工具——我们的手——的用途，它们轻易跟随每一种技艺和操作，无论在战争还是和平中都有用，但自然将它们加于身体，主要是为了理性。因为如果人缺少手，他面部的各部分肯定会像四足动物那样安排，以适合其进食的目的：所以面部形状会拉长，并指向鼻孔，嘴唇会从口中突出，块状、僵硬、厚实，适宜叼草；舌头要么躺在牙齿间，与嘴唇匹配，肉质、坚硬、粗糙，协助牙齿处理臼齿下的东西；要么潮湿，垂在一边，像狗及其他肉食动物，通过参差不齐的齿缝伸出。那么，如果我们的身体没有手，清晰的声音怎能被植入其中呢？因为口部各部分的形状就不会有适合言语使用的构造，那么人必然要么咩咩叫，要么\Quote{咩咩叫}，要么汪汪叫，要么嘶鸣，要么像牛驴一样吼叫，或发出某种野兽的声音？但现在，既然手成为身体的一部分，口便得闲暇服务于理性。由此可见，手是理性本性的特质，造物主藉此专门为理性设计了便利。\par

% structure: p0043 source=h2 target=section reason=relative-heading-rank; base=h2

\section{第九章。\Emph{人的形体被塑造为理性运用的工具}。}

1. 既然我们的造物主将某种\OriginalTerm{天主似的恩宠}{Godlike grace}赐予我们的形成，通过在自己本性的卓越形象上印上祂肖像的相似性，因此，祂慷慨地将其他美善恩赐赋予人性；但心智和理性（\OriginalTerm{心智}{mind}、\OriginalTerm{理性}{reason}）我们严格地说不能称祂\Emph{赐予}，而是\Emph{分施}，给肖像加上了祂本性的适当装饰。如今，既然心智是理智和非物质的，它的恩宠本是不可传达和孤立的，如果它的运动不通过某种机件显明出来。因此，仍然需要这种工具性的结构，使它像拨片一样，触动发声器官，并通过所发出的音符的性质，指示内在的运动。\par

2. 如同一位技艺熟练的乐师，他可能因某种疾病失去了自己的声音，但仍希望让人知道他的技艺，便借他人的声音奏出旋律，并借助长笛或竖琴展示他的艺术；同样，人类的心智，作为各种概念的发现者，由于它无法仅凭\OriginalTerm{灵魂}{soul}向那些通过身体感官聆听的人揭示其理解的运动，便像一位技艺精湛的作曲家，触动这些有生命的乐器，并通过它们所产生的声音，显明它隐藏的思想。\par

3. 人类乐器的音乐，是一种长笛与竖琴的复合体，如同在合奏曲中共同谐鸣。因为当说话者发出言语的冲动使和声成音时，气息从充满空气的容器中通过气管被迫上升，并撞击那些以圆形排列分割这条笛状通道的内部突起，便以一种方式模仿通过长笛发出的声音，被膜质突起环绕驱动。而上颚在自身的凹陷中从下方接收声音，通过两条延伸至鼻孔的通道，以及穿孔骨骼周围的软骨，犹如通过某种鳞片状的突起，使它的回响更加响亮；而脸颊、舌头、咽部的机制——通过它，下巴在收缩时放松，在伸展至一点时绷紧——所有这些都以多种方式回应拨片在弦上的运动，根据需要非常迅速地改变音调的排列；嘴唇的开合，则如同演奏者根据旋律的节拍用手指阻挡长笛气息所产生的效果。\par

% structure: p0047 source=h2 target=section reason=relative-heading-rank; base=h2

\section{十、\Emph{心智通过感官运作。}}

1. 既然心智通过我们工具性的结构产生理性的音乐，我们生来就是理性的；而我认为，如果我们不得不动用嘴唇来满足身体的需要——提供食物这一繁重而辛劳的任务——我们本不会拥有\OriginalTerm{理性的恩赐}{gift of reason}。然而，实际情况是，我们的双手接管了这一服务，使嘴巴空出来用于理性的服务。\par

2. 不过，这个乐器的作用是双重的：一是产生声音，二是从外部接受概念；并且这两种\OriginalTerm{官能}{faculty}并不相互混合，而是各自留在本性所指定的作用中，不干涉其相邻者——既不是听觉官能承担说话，也不是说话官能承担听觉；因为后者总是在发出声音，而耳朵，正如\Person{撒罗满}在某处所说，\Quote{耳朵听不足}。\par

3. 关于我们内在官能的那一点，在我看来甚至尤为令人惊奇的是：那容纳一切由听觉倾入之物的\OriginalTerm{内在容器}{inner receptacle}有多大？谁是在其中记录所带进来的话语的书记？存贮被听觉放入的概念的仓库是什么样的？而且，为什么当许多不同种类的概念相互挤压时，所存储的事物之间不会产生混乱和错误？人们也可能对视觉的作用产生同样的惊奇感；因为通过视觉，心智同样领会那些身体外在的事物，并将现象的形象吸引到自身，在自身中铭记所见之物的印象。\par

4. 正如有一座广阔的城市，通过不同的入口接待所有来者，但所有人并不会聚集在某个特定地点，而是有些人去市场，有些人去房屋，有些人去教堂，或街道、小巷、剧院，各人按自己的喜好——我似乎在我们内察觉到这样一座\OriginalTerm{心智之城}{city of our mind}，它被通过感官的不同入口不断充满，而心智区分并检查每一样进入的事物，将它们安置在各自的知识部门中。\par

5. 正如沿着城市的比喻，常常可能发生这样的事情：同一家族和亲属的人并不从同一个门进入，而是从不同的入口进入，但一旦进入城垣之内，他们又重新聚集，彼此亲密无间（也可能发生相反的情况：陌生且互不相识的人常常共用一个入口进入城市，但共同的入口并不将他们联系起来；因为即使在城内，他们也可以分开，去寻找自己的亲属）；我似乎在\OriginalTerm{心智的广阔领域}{spacious territory of our mind}中也察觉到类似的情况；因为常常我们从不同的感官官能所获得的知识是同一的，如同同一对象在感官上被分割成几个部分；而相反，我们又可能通过某个感官学到许多彼此毫无关联的、多样的事物。\par

例如——因为用实例阐明论证更好——让我们假设我们正在探究味觉的特性，即什么是甜的，什么是品尝者应避免的。我们通过经验发现胆汁的苦味和蜂蜜的甜味；但当这些事实被知晓后，这种知识是通过味觉、嗅觉、听觉、通常还有触觉和视觉（同一件事以多种方式引入我们的理解）赋予我们的。因为当一个人看到蜂蜜、听到它的名字、尝到它的味道、通过嗅觉辨识它的气味、通过触觉测试它时，他通过每一种感官认识同一件事。\par

另一方面，我们通过某一种感官获得多样而多形式的信息，因为听觉接收各种声音，我们的视觉通过观看不同种类的事物来运作——它同样触及黑色和白色，以及所有因颜色对立而区分的事物——同样，味觉、嗅觉、触觉也是如此；每一种知觉都通过自身的感知能力将各种事物的知识植入我们内。\par

% structure: p0055 source=h2 target=section reason=relative-heading-rank; base=h2

\section{XI. \Emph{论心智的本性是看不见的。}}

那么，这个将自己分配给感官能力、并通过每种感官妥善接受事物知识的心智，按其本性是什么呢？它异于感官，我想没有一个理性人会怀疑；因为如果它等同于感官，它就会将感官运作的特殊特性归为单一，理由是它自身是单纯的，而在单纯之物中找不到多样性。但现在，既然大家都同意触觉是一回事，嗅觉是另一回事，而其他感官彼此之间也处于同样无法相互交流或混合的境地，我们必定认为，既然心智恰当地存在于每种感官中，它必是异于感知本性的某物，这样就没有变异可能附着于可理解之物。\par

\Quote{有谁知晓了主的心意？\BibleRef{罗 11:34}}宗徒问道；我进一步问：有谁理解自己的心智？让那些认为天主本性在其理解之内的人告诉我们，他们是否理解自己——是否知道自己心智的本性。\Quote{它繁多而极其复合。}那么，可理解之物如何是复合的呢？或者，不同种类的事物混合的方式是什么？又或者，\Quote{它是单纯的、非复合的。}那么它如何分散到感官的众多划分中？如何在统一中存在差异？如何在差异中维持统一？\par

但我藉助于天主的圣言本身找到这些难题的解答；因为祂说：\Quote{让我们照我们的\OriginalTerm{肖像}{image}，按我们的\OriginalTerm{相似}{likeness}造人\BibleRef{创 1:26}}。肖像之所以是肖像，只有在其所显示的那些我们于原型中知觉的属性无一缺失的情况下才成立；但它在哪一点上偏离了与原型的相似，就在哪一点上不再是肖像；因此，既然我们在天主性内静观的一个属性是本质的不可理解性，显然在这点上肖像必须能够显示它对原型的模仿。\par

因为，如果原型超越理解，而肖像的本性却被理解，那么我们所静观的两者属性的相反特征就会证明肖像的缺陷；但由于我们的心智——造物主的相似——的本性逃避我们的知识，它便精确地相似于那更高的本性，以其自身的不可知性象征那不可理解的本性。\par

% structure: p0060 source=h2 target=section reason=relative-heading-rank; base=h2

\section{XII. \Emph{论\OriginalTerm{主导原则}{ruling principle}应被视为居于何处的问题；其中也讨论了眼泪与笑，以及关于物质、本性与心智之间相互关系的生理学思辨。}}

那么，让所有那些将理智活动局限于某些身体器官的虚妄与臆测讨论终结吧；他们中有些人将主导原则置于心脏，另一些人则说心智居于脑中，并用一些似是而非的表面理由强化这些意见。因为那位将主导权归给心脏的人，以心脏的位置（它似乎占据身体的中间位置）作为其论证的证据，理由是意志运动容易从中心分布到整个身体，继而付诸行动；他还将人烦恼和激情的性情作为其论证的见证，因为这些情感似乎使此部分产生同感。另一方面，那些将大脑奉献给推理的人说，头部被本性建造成整个身体的某种堡垒，心智像国王一样居住在其中，感官如同使者和持盾卫士环绕它。他们也在以下事实中找到其意见的标记：那些脑膜受损的人推理异常扭曲，而那些因醉酒而头重的人忽略了得体之事。\par

2. 凡支持这些观点的人，都提出了一些更具物理性质的理由来支持他们关于主宰官能的意见。一人声称，由理智产生的运动在某种方式上与火的本性相似，因为火和理智都处于永恒运动之中；而且由于热被认为源自心脏区域，他便据此断言，思想的运动与热的运动性相结合，并断言热所封闭的心脏是理性本性的容器。另一人则声称，脑膜（他们如此称呼环绕大脑的组织）可以说是所有感官的基础或根，从而证实了他自己的论点，理由是理智的能量不可能存在于其他部分，只有在耳朵与其相连并撞击所接收的声音之处，视觉（自然属于眼睛所在部位的空腔）通过落在瞳孔上的影像形成内在表象，气味的性质通过鼻子吸入而被辨别，味觉则通过脑膜的测试，该膜从自身通过颈椎向下向喉峡通道发出敏感神经过程，并与那里的肌肉结合。\par

3. 我承认，灵魂的理智部分常因情欲的泛滥而受扰乱；理智因身体的某种意外而变得迟钝，以致阻碍其自然运作，这是事实；而且心脏可以说是身体中火元素的一种源头，并随情欲的冲动而运动；此外，除了这一点，我不否认（因为我从那些从事解剖研究的人那里也听到非常相同的说法）脑膜（按照持这种生理学观点者的理论）包裹着大脑，并浸没在从大脑散发出的蒸汽中，形成感官的基础；但我并不以此作为证据，证明非物质的本质因任何空间的界限而受到限制。\par

4. 当然，我们知道精神错乱并不仅仅由头部沉重引起，而是有经验的医生宣称，我们的理智也会因侧面的膈膜患病而减弱，他们将此病称为狂乱，因为那些膈膜的名称是\OriginalTerm{膈}{\GreekText{φρένες}}。而由悲伤引起的感觉被错误地认为产生于心脏；因为疼痛的并非心脏，而是胃的入口，人们却无知地将这种感觉归因于心脏。然而，那些仔细研究过这些感觉的人给出了如下的解释：由于在悲伤状态下，全身毛孔自然地收缩和闭合，一切在通过时遇到阻碍的东西都被驱赶到身体内部的空腔中，因此（也因为呼吸器官受周围组织的压迫），呼吸常常在自然的努力下变得更猛烈，试图扩大收缩的部分，从而打开被压缩的通道；我们将这种呼吸视为悲伤的症状，称之为呻吟或尖叫。此外，似乎压迫心脏区域的那种痛苦感觉，并非心脏的疼痛，而是胃的入口的疼痛，并且出于同样的原因（我指的是毛孔的压缩），容纳胆汁的容器收缩，将那种苦涩而刺激的汁液倾注到胃的入口；其证据是，悲伤者的面色变得苍白并呈黄疸色，因为胆汁因过度压力将其汁液注入血管。\par

5. 此外，相反的情绪，我指的是欢乐和欢笑，也有助于确立这一论点；因为身体的毛孔，在那些因听到愉快之事而融化于欢乐的人身上，也以某种方式溶解和放松。正如在前一种情况下，毛孔细微而不可察觉的呼出因悲伤而受阻，并且当它们压缩上部内脏的内部结构时，将潮湿的蒸汽驱向头部和脑膜，这些蒸汽因被大脑空腔过度保留而通过其底部的毛孔被驱出，同时眼睑的闭合以水滴的形式排出湿气（那滴被称为眼泪），同样，我希望你理解，当毛孔因相反状态而异常扩大时，一些空气通过它们被吸入内部，然后又通过口腔的通道被自然排出，而所有内脏（尤其是他们说肝脏）通过某种激动和跳动的运动参与排出这种空气；由此，自然为了便于空气的排出，扩大了口腔的通道，将脸颊向两侧伸展在气息周围；其结果被称为欢笑。\par

6. 因此，我们既不应因此而将主宰官能归于肝脏，也不应因为愤怒情绪中心脏周围血液的发热状态而认为思想的座位在心脏；而必须将这些事情归于我们身体组织的特性，并认为思想以一种不可描述的联合方式同样与每个部分接触。\par

7. 即使有人在这方面引用圣经，声称心是主宰官能，我们也不应不加审视地接受这一说法；因为那提到心的人也提到肾，他说：\BibleQuote{天主考验心与肾}；因此他们要么将理智官能限于两者结合，要么限于两者都不。\par

虽然我意识到，理智的能力在身体的某种状态下会变得迟钝，甚至完全失效，但我并不认为这足以证明\OriginalTerm{理智的官能}{faculty of the mind}被限制在某个特定的地方，以至于因身体邻近部位可能出现的炎症而被挤出其应有的自由空间（因为这种看法是物质性的：当容器已被某种置于其中的东西占据时，其他东西便无法容身）；因为理智的本性既不住在身体的虚空中，也不会因肉体的侵占而被排挤；相反，整个身体被造成某种乐器，正如在那些懂得演奏的人身上常常发生的那样，他们无法展示技巧，因为乐器的不适不允许他们施展技艺（因为被时间毁坏、或因跌落而破损、或因生锈或腐朽而无用的乐器，即使被一位优秀的笛手吹奏，也是无声而无效的）；同样，理智贯穿整个乐器，以其理智活动相应的方式接触每个部分，按其本性，对那些处于自然状态的部分产生适当的效果，但在那些无法接纳其技艺运动的部分上则保持静止和无效；因为理智天生就与处于自然状态的事物密切关联，而与偏离自然状态的事物疏离。\par

在这里，我认为有一种更接近自然的看法，借此我们可以学到一些更精微的道理。既然万有中最美、至高的善就是\OriginalTerm{天主本身}{Divinity Itself}，一切倾向于\OriginalTerm{美善}{beautiful and good}的事物都朝向祂，我们因此说，理智作为最美者的肖像，只要它尽可能分有与\OriginalTerm{原型}{archetype}相似的特性，它自身就保持在美善之中；但如果它稍有偏离，它就被剥夺了所拥有的那种美。正如我们所说，理智因与原型之美的相似而受到装饰，仿佛被造成一面镜子来接受它所表达的原型形象；我们认为，由理智所治理的本性以同样的关系依附于理智，并且它也因理智所给予的美而受装饰，可以说是\OriginalTerm{镜中之镜}{mirror of the mirror}；而物质元素——即本性在其中被思考的那个存在——也受其支配和维持。\par

因此，只要一方与另一方保持接触，\OriginalTerm{真美}{true beauty}的交流就按比例贯穿整个系列，通过较高的本性美化其次者；但当这种有益的连接发生中断时，或者当较高的本性反而跟随较低的本性时，物质的无形丑态便显现出来——当它与本性分离时（因为物质本身是无形式、无结构的），由于它的无形，本性通过理智所装饰的那种美也被破坏；这样，物质的丑陋通过本性传递到理智本身，以致\OriginalTerm{天主的肖像}{image of God}不再见于那按肖像模塑之物的形象中；因为理智将美善的观念像镜子一样放在身后，转离了\OriginalTerm{至善}{the good}光辉的入射光线，反而接纳了物质的无形印记。\par

如此，邪恶的生成便是通过远离美善而发生的。凡是与\OriginalTerm{第一善}{First Good}密切相关的，都是美善的；但那些脱离与第一善的关系和相似的事物，当然缺乏美善。因此，按照我们正在讨论的说法，真正的善是唯一的，而理智自身也拥有美善的能力，只要它处于美善的肖像中；由理智所维持的本性也具有类似的能力，只要它是肖像的肖像；由此表明，我们物质的部分当被本性控制时便凝聚维持，否则当它与维持者分离、脱离与美善的结合时，便化解为无序状态。\par

这种状况除非在本性向相反状态翻转时才会发生，其中欲望不再倾向美善，而是倾向需要装饰的元素；因为那被塑造为与物质相似的东西，既然物质自身缺乏形式，必然也在缺乏形式与美这一点上与之同化。\par

然而，我们对这些点的讨论是随论证而进行的，因为它们是由我们对当前问题的思辨引出的；因为探究的主题是：\OriginalTerm{理智的官能}{intellectual faculty}是居于我们身体的某一部分，还是均匀地遍及全体？至于那些将理智局限于身体局部的人，他们为确立此见解而提出的事实是：那些脑膜处于非自然状态的人，理性无法自由运行。我们的论证表明：就人的复合本性的每一部分——由此每个人具有某种自然运作——而言，如果部分不保持自然状态，灵魂的能力同样保持无效。因此，沿着这条思路，我们刚才陈述的看法便进入我们的论证，由此我们得知：在人的复合本性中，理智受天主治理，而物质生活受理智治理，只要后者保持自然状态；但如果它偏离自然，它也就疏离于理智所进行的运作。\par

14. 然而，让我们再次回到我们的起点——即在那些没有被某种情感扭曲其自然状态的人身上，理智行使它自己的能力，并在那些健康的人身上坚定地建立；相反，在那些不接受其运作的人身上，它是无力的。因为我们还可以用其他论证来确证我们对这些事的意见：如果这对已经厌倦我们论述的听众不算是冗长的话，我们将尽我们所能，用几句话讨论这些事情。\par

% structure: p0075 source=h2 target=section reason=relative-heading-rank; base=h2

\section{XIII. \Emph{睡眠、打哈欠与梦幻的原理}}

1. 我们身体的这种生命，是物质且易变的，总是通过运动前进，它将其存在的能力维系于此，即它从不停止其运动；如同一条河流，凭借自身动力奔流，使河床保持满溢，但同样的水不会始终在同一位置，一部分流走，另一部分涌来，同样，我们这里生命的物质元素也在对立面的连续更替中通过运动和流动经历变化，以至于它永远无法停止变化，而是因其无法静止，通过类似的方式持续交替运动：如果它一旦停止运动，它必然也会停止存在。\par

2. 例如，充满之后是空虚，另一方面，空虚之后又轮流是充实的过程：睡眠放松清醒的紧张，而清醒又使松弛的变得紧张：这两者都不会持续存在，而是在对方到来时彼此退让；自然通过它们的交替如此更新自身，轮流参与两者，不间断地从一方过渡到另一方。因为生物若持续努力操作，会导致过度紧张部分的某种断裂和分离；而身体持续静止会导致结构的一定分解和松弛：但在适当的时候适度接触两者，是自然的保持力，它通过不断转移到对立状态，在每种状态中从另一方得到休息。因此，她发现身体因清醒而紧张，便通过睡眠来缓解紧张，暂时让感知能力从操作中休息，像比赛后把马从战车上解下。\par

3. 进一步，身体框架需要在适当时间休息，以便营养通过其所含的通道遍布全身，而没有任何紧张阻碍其进程。因为正如雨水浸泡大地后，当太阳用相当温暖的光线加热时，从地底升起某种雾状蒸汽，同样，在我们体内的土地中，当内部的营养被自然温暖加热时，也会发生类似的结果；蒸汽自然有上升趋势和空气性质，渴望居于其上者，便来到头部区域，像烟雾穿透墙壁的接缝；然后它们经呼气从那里散发到感觉器官的通道，感官自然因此变得不活跃，让位于这些蒸汽的通过。眼睛被眼皮按压，当某种像铅一样的工具（我是指像我说过的那种重量）将眼皮放低到眼睛上时；听力被同样的蒸汽弄钝，如同在听觉器官上安了一扇门，便从其自然操作中休息：这样的状态就是睡眠，此时感觉在身体中休息，完全停止其自然运动的操作，以便营养的消化过程可以通过蒸汽自由地通过每个通道。\par

4. 因此，如果感觉器官的装置被关闭，且睡眠因某种活动而受阻，神经系统充满蒸汽，自然自发地伸展，以至于因蒸汽而密度增加的部分通过伸展过程变得稀薄，就像那些通过用力拧干衣服中的水的人所做的一样；而且，由于咽喉周围的部分有些环状，且神经组织丰富，每当需要从该部分排出浓密的蒸汽时——因为不可能直接分开圆形形状的部分，只能通过其圆周轮廓的扩张——因此，通过打哈欠屏住呼吸，下巴向下移动，以便在小舌处留下一个凹陷，所有内部部分配置成圆形，那些困留在附近部分的烟雾状密度随呼吸排出而一同排出。即使在睡眠后，如果任何一部分蒸汽仍留在所说区域未被消化和呼出，也常会发生类似情况。\par

5. 因此，人的理智清楚地证明了其与自身本性相连的主张，理智本身也在健康清醒状态下与本性合作并一同运动，但在陷入睡眠时保持不动，除非有人认为梦幻的想象是理智在睡眠中的运动。我们却说，只有理智有意识和健康的行动才应归于理智；至于睡眠时发生的荒谬幻想，我们认为是一些理智操作的幻象偶然在灵魂中较不理性的部分中形成；因为灵魂通过睡眠与感官分离，也必然超出理智操作的范围；因为正是通过感官，理智与人的结合才发生；因此，当感官休息时，理智也必然不活动；证据是做梦者常处于荒谬且不可能的情境中，如果灵魂当时由理性和理智引导，就不会发生这种情况。\par

6. 然而在我看来，当灵魂就其更卓越的官能（我指的是心智和感官的活动）而言处于休息状态时，睡眠中只有其营养部分在运作；而我们在清醒时发生的事物——感官和理智的活动——的一些阴影和回响，被灵魂中能记忆的部分印刻在营养部分上，这些，我说，是随机会而显现的，记忆的一些回响仍萦绕在灵魂的这一部分。\par

7. 因此，人就被这些迷惑了，并非通过任何思想序列而被引导认识呈现的事物，而是在混乱无绪的错觉中游荡。但正如在其身体活动中，每个部分都根据其自然具有的能力以某种方式运作，在静止的肢体中也会产生一种与运动中的部分相共鸣的状态；同样，就灵魂而言，即使一部分静止而另一部分运动，整个灵魂也与那部分共鸣；因为自然的统一性不可能以任何方式被割裂，尽管其中所包含的官能之一因其活跃运作而依次占主导。但如同人们清醒忙碌时，心智占主导，感官服务它，而调节身体的官能并不与它们分离（因为心智提供其所需的食物，感官接受所供应的，身体的营养官能将给予的东西据为己有），同样，在睡眠中，这些官能的支配权在我们内以某种方式颠倒，当较不理性的部分成为主导时，另一官能的运作确实停止，却并未完全熄灭；但当营养官能那时在睡眠中忙于消化，使我们全部本性专注于自身时，感官官能既不完全与它分离（因为一旦自然结合就无法分离），也无法恢复其活动，因为感官器官在睡眠中不活动而受阻；按同样推理（心智也与灵魂的感官部分结合），我们应当说，心智在感官运动时随之运动，在感官静止时随之静止。\par

8. 正如火被谷壳覆盖，没有气息吹动火焰时，它既不会烧毁旁边的物体，也不会完全熄灭，而是以烟的形式通过谷壳升到空中；然而若它获得任何气息，便将烟变为火焰——同样，心智在睡眠中因感官不活动而被隐藏时，既不能通过它们闪耀出来，也不会完全熄灭，而是有一种，可以说，阴燃的活动，在一定程度上运作，却无法进一步运作。\par

9. 再者，正如一位音乐家，当他用拨子弹拨里拉琴松弛的琴弦时，奏不出有序的旋律（因为未绷紧的弦不会发声），但他的手常技巧地移动，将拨子放在音符的位置上，就位置而言，却没有声音，只是通过琴弦的振动产生一种不确定而模糊的哼声；同样，在睡眠中，感官的机制松弛了，艺术家要么完全不活动（如果琴弦因饱足或沉重而完全松弛），要么松弛而微弱地活动，如果感官的乐器不完全容许其艺术运作的话。\par

10. 因此，记忆是混乱的，而预知，虽被不确定的帷幕弄得模糊，却以我们清醒活动的阴影成像，并常常向我们指示将要发生的一些事：因为心智凭借其本性的精微，在洞察事物方面比单纯的肉体粗重有一些优势；但它无法通过直接方法使它的意思清楚，使所处理之事的信息明白而清晰，它对未来的宣告却是模糊而可疑的——那些解释此类事物的人称之为\Quote{\OriginalTerm{谜语}{enigma}}。\par

11. 所以司酒者挤压那串葡萄为\Person{法老}的杯；所以司膳者似乎携带他的篮子；每个人在睡眠中想象自己从事着清醒时忙碌的那些职务：因为他们习惯性职务的影像印在他们灵魂中有预知能力的部分，藉着心智的这种预言，暂时给予他们预知将要发生之事的能力。\par

12. 但如果\Person{达尼尔}和\Person{若瑟}以及类似他们的人，藉\OriginalTerm{神圣}{Divine}的能力，毫无知觉混乱地接受了未来事物的知识，这并不影响目前的陈述；因为没有人会将此归于梦的能力，否则他必然也会认为那些在清醒中发生的\OriginalTerm{神圣}{Divine}显现也不是奇迹般的异象，而是自然自发产生的结果。正如所有的人都由自己的心智引导，却有少数人蒙受了明显的神圣启示；同样，虽然梦的幻想以相似而同等的方式自然发生在所有人身上，但有些人（而非所有人）藉着他们的梦分享了某种更神圣的显示：而对于其余所有人，即使通过梦确实发生了对任何事物的预知，那也是按我们所说方式发生的。\par

13. 再者，如果\Place{埃及}王和\Place{亚述}王被天主引导认识未来，那么通过他们实现的安排是另一回事：因为有必要使圣人们隐藏的智慧得以显明，使他们每个人都不致对国家毫无益处而虚度一生。因为如果占卜者和术士不能胜任发现那梦的任务，\Person{达尼尔}又如何能因其所是而被认识？而\Person{若瑟}被关在狱中时，若非他解梦引起注意，\Place{埃及}又如何能得保存？因此我们必须将这些情况视为例外，而不将它们与普通的梦归为一类。\par

14. 但这种普通的做梦现象是所有人共有的，它以不同的方式和形式产生于我们的幻想中：因为要么是我们所说的日常事务的回声仍留在灵魂的记忆部分；要么，如常发生的，梦的构成是根据身体这样或那样的状况而形成的：例如，口渴的人似乎身处泉源之间，需要食物的人似乎身处宴席，而年轻人在青春活力的热力中被与他的情欲相应的幻象所困扰。\par

15. 我也知道梦的幻象的另一种原因，是在我照看一位被狂乱症侵袭的亲戚时发现的。他因被给予超其体力的过多食物而恼怒，不停地喊叫，指责身边的人往肠子里填满粪便并放在他身上；当他的身体迅速趋向出汗时，他责怪身边的人准备好了水，待他躺着时弄湿他。他不停地喊叫，直到结果显明了这些抱怨的含义：因为突然间全身大汗淋漓，而肠道的松弛解释了肠中的重量。那么，同样的情况——即在他清醒的判断因疾病而迟钝时，他的本性所经历的，与身体状况同感相应，并非没有察觉到不对劲，但由于疾病引起的心神涣散而无法清楚表达痛苦——这种情况，如果灵魂的理智原则被平息入眠，不是由于虚弱而是由于自然的睡眠，可能会作为一个梦出现在处境相似的人面前，汗水的爆发被表现为水，食物引起的痛苦被表现为肠子的重量。\par

16. 这个观点也为那些精通医学的人所采纳，即根据病情的不同，病人的梦的幻象也以不同的方式出现：胃弱的人有一种幻象，脑膜受损的人有另一种，发烧的人又有另一种；胆汁性和痰湿性患者的幻象不同，多血症和消瘦患者的幻象也不同；由此我们可以看出，灵魂的营养和生长功能中混合着某些理智元素的种子，这些种子在某种意义上与身体的特定状态相似，根据侵袭它的症状而在其幻象中适应。\par

17. 此外，大多数人的梦与他们的性格状态相应：勇士的幻象是一种，懦夫的又是另一种；放荡者的梦是一种，节制者的又是另一种；慷慨者和吝啬者有不同的幻象；而这些幻象绝非由理智形成，而是由灵魂的较不理性的倾向形成，这倾向甚至在梦中塑造出各人清醒时习以为常的事物的相似物。\par

% structure: p0093 source=h2 target=section reason=relative-heading-rank; base=h2

\section{XIV. \Emph{理智不在身体的某个部分；其中也区分身体和灵魂的运动}。}

1. 但我们偏离主题很远，因为我们论证的目的是要表明理智并不局限于身体的任何部分，而是同等与整个身体接触，根据其所影响的部分的性质产生运动。然而，有时理智甚至跟随身体的冲动，成为其仆人；因为身体的本性常常先导，引入要么痛苦的感觉，要么快乐的欲望，以至于可以说它提供了最初的开始，在我们心中产生对食物的欲望，或者一般来说，对某种快乐事物的冲动；而理智接收这种冲动，通过自身的智慧为身体提供达到所欲之物的恰当手段。确实，这种情况并非发生在所有人身上，只发生在那些较为奴性的人身上，他们将理性束缚在本性冲动之下，通过允许理性与之结盟而奴性地尊崇感官快乐；但在更成全的人身上，这种情况不会发生；因为理智主导，通过理性而非激情选择有益之举，而他们的本性则跟随其引领者的足迹。\par

2. 但既然我们的论证在生命官能中发现了三种不同的种类——一种无感觉地接收营养，另一种既接收营养又能感觉，但没有理性活动，第三种是理性的、完美的，并与整个官能并存——因此在这些种类中，理智的官能拥有优势；——但愿没有人因此认为在人的复合本性中有三个灵魂焊接在一起，各自在自己的界限内被默观，以致有人以为人的本性是若干灵魂的某种混杂。真正的、完美的灵魂本质上是唯一的，即理智的和非物质的，它通过感官与我们的物质本性结合；而一切物质本性，易变和可变的，如果参与赋生能力，就会藉生长运动；相反地，如果它从生命活力中坠落，就会将其运动导向毁灭。\par

3. 因此，没有物质实体就没有感觉，没有理智官能也没有感觉活动。\par

% structure: p0097 source=h2 target=section reason=relative-heading-rank; base=h2

\section{XV. \Emph{灵魂本身，在事实和名称上，是理性的灵魂，而其他的只是同名地被称为灵魂；其中也包含这一陈述：理智的能力通过适当的接触遍及整个身体的每一部分}。}

1. 如今，如果受造物中有些拥有滋养官能，另一些则由感知官能支配，前者没有感知，后者没有理智本性，若有人因此倾向灵魂多元论，此人将以不符合其区别定义的方式提出多种灵魂。因为我们在现存事物中所构想的一切，如果它完全就是其所是，那么它恰当地被称为它所具有的名称；但若它所称为的并非完全如此，则该称谓也是虚妄的。例如：若有人向我们展示真正的面包，我们说此人正确地将名称应用于该事物；但若有人展示的是用石头仿制天然面包的东西，它形状相同、大小相等、颜色相似，以至于在多数方面与原型相同，但却缺乏作为食物的能力，因此我们说这块石头被称为\Quote{面包}并不恰当，而是误称，所有属于同样描述的事物，若非绝对地就是它们所被称为的，它们的名称都源于术语的误用。\par

2. 因此，正如灵魂在理智和理性的事物中找到其完美，凡不是这样的事物，虽然可以分享\Quote{灵魂}之名，却并非真正的灵魂，而是与\Quote{灵魂}之称相关联的某种生命力。为此，那位对一切事物制定法律者，也将动物本性（这种本性并没有远离植物生命）赐予人使用，作为那些享用者代替蔬菜的食物——因为祂说：\Quote{你们要吃各种肉类，如同青菜一样；}因为感知官能似乎只比无需感知而滋养生长的东西略胜一筹。愿这教导属肉性的人不要将自己的理智紧紧束缚于感官现象，而要忙于精神方面的益处，因为真正的灵魂在这些里面，而感官在野兽中也具有同等的力量。\par

3. 不过，我们的论证进程已转向另一点：因为我们思考的主题并非在于，在我们所构想的人的各种属性中，心灵的活动比其物质成分更具尊贵，而在于心灵并不局限于我们身体的任何一部分，而是同等地在一切部分中并通过一切部分，既不包围任何外在的东西，也不被任何东西所包围：因为这些说法恰当地用于桶或其他一个套一个的物体；但心灵与身体的结合呈现一种不可言说、不可思议的联系——既不在其\Emph{内}（因为非物质的不能被身体包围），也不在其外包围它（因为非物质的不包含任何东西），而是心灵以某种无法言说、无法理解的方式接近我们的本性，并与本性接触，它应被视为既在其中又在其周围，既非植根于它，也非被它包裹，而是以我们无法言说或思考的方式，除了这一点：当本性按其自身秩序繁荣时，心灵也在运作；但如果前者遭遇任何不幸，智力的运动也随之停止。\par

% structure: p0101 source=h2 target=section reason=relative-heading-rank; base=h2

\section{\Emph{第十六章  默想天主圣言所说的——\Quote{让我们按我们的肖像和模样造人}；其中探讨肖像的定义，以及可苦的、可朽坏的如何相似于有福的、不可苦的，并且肖像中有男有女，而原型中没有这些。}}

1. 现在让我们重新考虑天主的圣言：\Quote{让我们按我们的肖像，按我们的模样造人。}\BibleRef{创 1:26} 一些异教作家的幻想是多么低下、多么与人的尊严不相称啊！他们通过将人性与这世界相比来抬高人性，如他们所想的那样。因为他们说人是小宇宙，由与宇宙相同的元素组成。那些以如此高调的名称将这种赞美赋予人性的人，忘记他们是用蚊虫和老鼠的属性来荣耀人：因为蚊虫和老鼠也由这四种元素组成——因为可以肯定，在每个存在物的有生命本性中，我们都看到那些元素的一部分或大或小，没有它们，任何有感觉的存在都不可能存在。那么，如果人被认为是世界的代表和肖像——即将逝去的天、变化的地、以及其中包含的一切，这些随着环绕它们之物的消逝而消逝——有什么伟大的呢？\par

2. 那么，按照教会的道理，人的伟大在于什么呢？不在于他与受造世界的相似，而在于他是造物主本性的肖像。\par

3. 那么，也许你会问，肖像的定义是什么？非物质的如何相似于身体？暂时的如何相似于永恒的？因变化而易变的那位如何相似于不变的？受激情和腐朽支配的那位如何相似于不可苦的、不朽的？常与邪恶同住、在邪恶中成长的那位如何相似于完全免于邪恶的？因为在原型中所构想的事物与按其肖像所造的事物之间有很大的区别：因为肖像之所以恰当地称为肖像，就在于它保持与原型的相似；但如果模仿偏离了其主体，那事物就成了别的东西，而不再是主体的肖像。\par

4. 那么，这个可朽的、可受难的、短命的人，如何成为那不朽的、纯洁的、永恒之本性的肖像呢？对于这个问题，或许只有真理本身才知道真正的答案；但这就是我们通过猜测和推理，尽可能追溯真理所领悟到的。天主的言语并没有说谎，它说人是按天主的肖像造的；人的本性之可悲苦难也不像那不受难生命的福乐；因为若有人将我们的本性比作天主，必然要承认两件事之一，才能使相似的定义在两种情况下以相同的术语被理解——要么天主神性是可受难的，要么人性是不受难的；但若天主神性既非可受难，我们的本性也并非无苦难，那么还有什么说法能让我们说天主的言语是真实的，它说人是按天主的肖像造的呢？\par

5. 那么，我们必须再次拿起圣经本身，也许我们可以从所写的文字中找到一些指引。在说了\Quote{让我们按我们的肖像造人}以及为什么说\Quote{让我们造他}之后，它又加上了这句话：——\Quote{天主于是造了人；按天主的肖像造了他；造了他们男和女\BibleRef{创 1:27}。}我们在前面已经说过，这句话是为了摧毁异端的不敬而说的，为使我们得知独生天主按天主的肖像造了人，就绝不将父和子的天主性加以区分，因为圣经对造人的那一位和按祂肖像所造的那一位，都同样赋予天主的名称。\par

6. 然而，让我们略过关于这点的论证；我们把探究转向眼前的问题——既然天主神性在福乐中，而人性在苦难中，后者为何仍在圣经中被称为与前者的\Quote{相似}？\par

7. 那么，我们必须仔细审查这些话：因为若如此行，我们就会发现那按\Quote{肖像}所造的是一回事，而如今显现在悲惨中的是另一回事。经上说：\Quote{天主造了人，按天主的肖像造了他\BibleRef{创 1:27}。}那按\Quote{肖像}所造的被造物到此为止；然后它又重拾创造的叙述，说：\Quote{造了他们男和女。}我想每个人都知道这是与原型的偏离：因为如宗徒所说，\Quote{在基督耶稣内，既不是男也不是女。}然而这句话却宣称人是如此分裂的。\par

8. 因此，我们本性的创造在某种意义上可以说是双重的：一个是像天主的，另一个是按这种区分而分开的；因为这段经文通过它的安排隐约传达了这样的意思，它首先说：\Quote{天主造了人，按天主的肖像造了他\BibleRef{创 1:27}，}然后补充道：\Quote{造了他们男和女\BibleRef{创 1:27}——}这是与我们关于天主的观念格格不入的。\par

9. 我认为圣经通过这些话语向我们传达了一个伟大而崇高的道理；这个道理是这样的。虽然有两种本性——神圣的非物质本性和非理性的畜生生命——作为极端彼此分离，但人的本性是它们之间的中道：因为在人的复合本性中，我们可以看见我所提及的两种本性的各部分——神圣方面的，是理性与理智的元素，它不接受男女的区分；非理性方面的，是我们的身体形态与结构，分为男和女：因为每个元素确实存在于一切享有人的生命者之中。然而，理智的元素先于另一个，我们从一位按顺序叙述人的创造者那里得知；我们也得知他与非理性的共同之处和亲缘关系，是为人的生育而预备的。因为他首先说\Quote{天主按天主的肖像造了人}（用这些话表明，如宗徒所说，在这样的存在中没有男或女）；然后他加上人性的特性：\Quote{造了他们男和女\BibleRef{创 1:27}。}\par

10. 那么，我们从中学到了什么？请任何人不要愤怒，如果我从远处引一个论证来支持眼前的话题。天主按其本性就是我们的心所能想到的一切美善——不如说，超越我们所能思想或领悟的一切美善。祂造人没有别的理由，只因为祂是美善的；而祂既是如此，并且以此为进入创造我们本性的理由，祂就不会以不完美的形式显示祂美善的能力，只给予我们本性祂所拥有的一些东西，而吝啬分享另一种；但美善的完美形式在此可见，是祂既从虚无中带人进入存在，又用一切美善的恩赐充分供应他；但由于个别美善恩赐的清单很长，无法用数目来把握。因此圣经的语言通过一个概括性的短语简洁地表达出来，说人是按\Quote{天主的肖像}造的：因为这就等于说祂使人性分有了一切美善；因为若天主神性是美善的圆满，而这是祂的肖像，那么这肖像在充满一切美善上就找到了与原型相似之处。\par

11. 因此，在我们内有了一切卓越、一切德性与智慧，以及我们所思想的每一种更高的事物；但其中最卓越的是我们不受必然性束缚，不受任何自然力量的奴役，而有权按我们的意愿自行决定；因为德性是一种自愿的事，不受任何辖制；强迫和暴力所导致的结果不可能是德性。\par

12. \OriginalTerm{肖像}{image}在各方面都带有\OriginalTerm{原型}{archetype}的卓越的相似，如果它在某方面没有差异，那么完全无分歧就不再是相似，而显然与\OriginalTerm{原始模型}{Prototype}完全相同。那么我们在神性与那被造成相似于神性者之间发现什么差异呢？我们发现前者是\OriginalTerm{非受造}{uncreate}的，而后者是由受造而存在；这个特性的区别带来了一系列其他特性；因为非常确定地承认，非受造的本性也是\OriginalTerm{不变的}{immutable}的，始终如一，而受造的本性不能没有变化而存在；因为它从不存在到存在的过渡本身，就是非存在之物被天主的旨意转化为存在的一种\OriginalTerm{运动}{motion}与\OriginalTerm{变化}{change}。\par

13. 正如福音称钱币上的印记为\Quote{凯撒的肖像}\BibleRef{玛 22:20}，由此我们得知，在那被塑造得相似于凯撒的肖像上，有外貌的\OriginalTerm{相似}{likeness}，而\OriginalTerm{材料}{material}不同，同样，在当前的论述中，当我们考虑那构成相似性、存在于神性与人性两者中的属性时，这些属性相当于面貌特征，我们在那作为基础的东西中发现了我们在非受造的本性与受造的本性中所见的差异。\par

14. 正如前者始终如一，而那由受造而存在的起始就源于变化，并与类似的变化有着同源的联系，为此，那如先知著作所说\Quote{预先知道一切将成之事}者，借着祂预知的能力，推演，或更确切地说，预先感知人在独立和自由状态下其意志的趋向如何，——因为祂预见了将要发生之事，祂为祂的肖像设计了男和女的区别，这与\OriginalTerm{神圣原型}{Divine Archetype}无关，而是如我们所说，是更接近\OriginalTerm{较不理性的本性}{less rational nature}。\par

15. 实际上，这一设计的原因，只有那些作为真理的目击者和\OriginalTerm{圣言}{Word}的仆役的人才能知道；而我们凭借推测和比喻，尽可能设想真理，并不以权威的方式陈述我们心中所想，而是以理论推测的形式呈现在我们友善的听众面前。\par

16. 那么我们关于这些事理解什么呢？经文说\Quote{天主创造了人}，用不明确的术语表示全人类；因为亚当不是与创造一起命名的吗？正如历史在随后的叙述中告诉我们？然而，被创造的人的名称不是特指的名称，而是总称：因此，通过使用我们本性的总称，我们被引导到这样的看法——在天主的预知和能力中，全人类都包含在第一次创造中；因为天主看祂所造的一切都不是不确定的，而是每一存在之物都有其创造者智慧所规定的一些界限和尺度，这是适宜的。\par

17. 正如任何一个特定的人都被他的身体尺寸所限制，与身体表面相结合的特定大小是他个别存在的尺度，同样，我认为全人类的整个圆满都被万有的天主借着预知的能力包含在如同一个身体中，这就是经文教导我们的，它说\Quote{天主创造了人，按天主的肖像造了他}。因为\OriginalTerm{肖像}{image}并不在于我们本性的某一部分，也不在于本性中某一单独的事物上，而是这能力平等地延伸到整个种族：这标志是理智同样植入在所有人中：因为所有人都有理解和思辨的能力，以及其他一切使神性本体得以在依它而造者身上显现其肖像的能力：在世界最初创造时显现的人，和在万物终结之后的人，都是一样的：他们同等地在自己身上带着天主的肖像。\par

18. 为此，整个种族被说成是一个人，就是对于天主的权能而言，没有过去或未来，甚至我们所期待的事物，也同现存的事物一样，被支持万有的能量所包含。我们的整个本性，从最初到最终，可以说是\OriginalTerm{自有者的肖像}{image of Him Who is}；但男女性别的区分是在祂造物工作的最后才添加的，我认为，是出于如下的原因。\par

% structure: p0120 source=h2 target=section reason=relative-heading-rank; base=h2

\section{XVII. 应当如何回答那些提出此问题的人——\Quote{如果生育在罪恶之后，那么如果人类的第一对夫妇仍保持无罪，灵魂将如何产生？}}

1. 然而，对我们来说，在探究这一点之前，或许最好先探讨一下我们的对手提出的问题的解答；因为他们声称，在罪恶之前，没有记载关于\OriginalTerm{生育}{procreation}，或\OriginalTerm{分娩之痛}{travail}，或趋向生育的欲望；但当他们被逐出乐园之后，女人因分娩之痛而被判罪，亚当这才与他的伴侣进入婚姻生活的交往，然后才有了生育的开始。那么，如果乐园中没有\OriginalTerm{婚姻}{marriage}，没有分娩之痛，没有生育，他们声称，必然的结论是，如果那\OriginalTerm{不死}{immortality}之恩没有堕落到\OriginalTerm{死亡}{mortality}，并且婚姻没有通过后代来保存我们的种族，用新生来代替离去者，那么人类的灵魂就不会有众多存在，这样在某种程度上，进入世界的罪对人类生命是有益的：因为人类种族本会停留在第一对夫妇中，要不是\OriginalTerm{对死亡的恐惧}{fear of death}驱使他们的本性去提供延续。\par

2. 现在，真正的答案——不论它是什么——唯有那像\Person{保禄}一样，曾受教于乐园奥迹的人才能明了；但我们的答案如下。当撒杜塞人曾为反对复活道理而辩论，并提出那曾有七位兄弟为夫的多次结婚的妇人来支持他们自己的意见，问她复活后将是哪一位的妻子时，我们的主回答他们的论点，不仅教训了撒杜塞人，也启示了后世所有人有关复活生命的奥迹。祂说：\BibleQuote{因为在复活时，他们不娶也不嫁；他们也不能再死，因为他们与天使相似，既是复活之子，也就是天主的子女。\BibleRef{路 20:35}}。现在，复活应许给我们的，不外乎是使堕落者恢复其原始状态；因为我们所期待的恩宠，就是重返最初的生命，使那被逐出乐园者再次回归乐园。既然那些得恢复者的生命与天使的生命密切相关，那么显然，犯罪之前的生命是一种天使般的生命，因此我们生命的原始状态的回归也被比拟为天使。然而，如前所述，在他们中没有婚姻，天使的军队却有无数万；因为\Person{达尼尔}在他的神视中这样宣告。同样，如果不是因罪恶导致我们堕落、从与天使平等的地位被移除，我们也不必需要婚姻来繁衍；但不论天使本性的增长方式如何（不可言说、也不可凭人的猜想理解，除非它确实存在），它也曾在人身上运作，人原是被造得\Quote{稍微逊于天使}，为要使人类增长至造物主所定的数目。\par

3. 但如果有人对灵魂生成的方式感到困难——若人不需要婚姻的帮助，我们将反问他们：天使的存在方式如何？他们如何以同一本质却同时有无数万之多？因为我们将给出一个恰当的答案给那提出\Quote{人如果没有婚姻会怎样}问题的人，如果我们说：\Quote{正如天使没有婚姻一样}；因为人在犯罪之前与他们处于相似的状态这一事实，由重返那状态而显明。\par

4. 既然我们已经澄清了这些事，让我们回到先前的论点——在造了祂的肖像之后，天主如何为祂的工程安排了男女之别。我说，我们已完成的初步思辨有助于确定这个问题；因为那创造万有、并凭自己的旨意将人整体地塑造为天主的肖像的那一位，并没有等待灵魂的数目经由后来者的逐渐补充而达到其适当的满盈；而是凭着祂的预知，观看人的本性在整体中的完满，并赐予它一份崇高而与天使相等的命运，因为祂以全知预见了他们意志对正直追求美善的失败，以及由此导致的与天使生命分离；为了不使人类灵魂的众多因从天使增长繁衍的方式坠落而减少——为此，我强调，祂为我们的本性形成了那适合于已陷于罪恶者的增长机制，将动物性的、无理性的方式植入人类，以取代本性中天使般的尊贵，人类现在就是以这种方式相继繁衍。\par

5. 因此，在我看来，伟大的\Person{达味}怜悯人的痛苦，用这样的话哀叹人的本性：\Quote{人虽有\InnerQuote{尊荣}却不自知}（以\Quote{尊荣}指与天使平等），因此，他说：\Quote{他被人比作无知的畜生，成为跟它们一样的}。因为那人真正变得像畜生一样，他在本性中接受了现今这种暂存繁衍的方式，出于他倾向于物质事物。\par

% structure: p0126 source=h2 target=section reason=relative-heading-rank; base=h2

\section{XVIII. \Emph{论非理性的情欲源自与无理性本性的亲缘。}}

1. 因为我认为，所有的情欲都从这初始像泉水一样涌出，泛滥于人的生命；我这话的证据就是情欲的亲缘关系，它同样出现在我们自己身上和禽兽身上；因为将我们本性上易于动情的最初根源归于那按天主肖像塑造的人性，是不允许的；但既然禽兽的生命首先进入世界，而人，因先前所述的理由，取了它们本性中的某些东西（我指生养模式），他同时也分享了那本性中可观察的其他属性；因为人与天主的相像并不见于愤怒，快感也不是那高超本性的标记；怯懦、鲁莽、贪得、厌失，以及一切类似的情感，都远非那显示神性的印记。\par

2. 这样，人性从禽兽方面取了这些属性；因为禽兽生命为自保而装备的那些素质，当转移到人的生命中时，就成了情欲；因为肉食动物靠其愤怒而存活，大量繁殖的动物靠其对快感的喜爱；怯懦保护弱者，恐惧保护易被更强大动物捕捉者，贪婪保护大块头者；而错过任何带来快感的东西，对禽兽来说是痛苦的事。所有这一切以及类似的情感，因动物的生养模式而进入人的构造。\par

3. 容我借某种奇妙的塑像来比拟人的形象。就如人们在模型中看到那些工匠为令观者惊奇而雕琢出的形状，在同一头上刻出两副面容；同样，在我看来，人也似乎同时肖似两种对立之物——在心灵的属神部分被塑造为神圣的美，却在自身涌起的激情冲动中肖似兽性；甚至其理性也常常沦为兽性，因趋向非理性之物而偏斜和倾向，以恶者遮蔽了善者；因为每当人将心智的力量拖向这些情感，并迫使理性成为激情的奴仆，那良善的印记便在某种程度上转化为非理性的形像，其整个性情便按那样式重新勾勒，理性可说是在培育激情的开端，并逐渐使之倍增；因为理性一旦与激情合作，就会产生丰盛而众多的恶果。\par

4. 因此，我们的爱乐起于与无理性受造物的相似，并由人的过犯增加，成为如此之多从快乐生出的罪愆之母，为无理性动物中所未见。同样，我们心中的愤怒确实与野兽的冲动相近，但因思想的联盟而增长：从此生出恶意、嫉妒、欺诈、阴谋、伪善；这一切都是心智不良耕作的结果；因为若激情被剥夺了从思想获得的协助，余下的愤怒便短暂而不持久，如同泡沫，刚一生成便立即消散。同样，猪的贪婪引入了贪财，马的高傲成为骄傲的根源；所有从无理性兽性中产生的特殊形态，因心智的恶用而成为恶行。\par

5. 同样，反之，如果理性反而掌控了这类情感，每一种都会被转化为美德的形态；因为愤怒产生勇气，恐惧产生谨慎，畏惧产生服从，憎恨产生对恶的厌恶，爱的力量产生对真正美善的渴望；性情中的高昂使我们的思想超越激情，免于受卑劣之事的束缚；是的，那位伟大的宗徒甚至称赞这种精神提升的形态，他嘱咐我们要常常\Quote{思念上面的事\BibleRef{哥 3:2}}；如此我们便发现，每一种这样的动因，一旦被思想的高尚所提升，便与神圣形像的美相符。\par

6. 但另一种冲动更大，因为罪的倾向沉重而向下；我们灵魂的统治要素更容易被无理性本性的重量向下拖曳，而非沉重和属土的因素被理智的高昂所提升；因此包围我们的痛苦常使神圣恩赐被遗忘，并将肉体的情欲如同丑陋的面具覆盖在形像的美上。\par

7. 因此，那些观看这类情形而不承认神圣形像存在的人，在某种意义上情有可原；然而，通过那些已妥当安排自己生活的人，我们可以在他们身上看见神圣的形像。因为，如果受情欲和属血气之人使人难以相信人曾被装饰以神圣之美，那么品德高尚、纯洁无染的人必会证实你对人性更美好的看法。\par

8. 例如（因为用实例使论证更清晰更好），某个以邪恶著称的人——比如某个耶苛尼雅，或其他恶名昭彰者——以邪恶的污染抹去了自己本性的美；然而在梅瑟和类似他的人身上，形像的样式保持纯洁。如今，当形像的美没有被遮蔽时，那句说人是天主的形像的话语的真实性便显明了。\par

9. 然而，或许有人因我们的生命像禽兽一样靠食物维持而感到羞愧，并因此认为人不配被假定为按天主的形像被造；但他可能期望，在我们所盼望的生命中，这一功能终将从我们的本性中被除去；因为正如宗徒所说：\Quote{天主的国不在于吃喝\BibleRef{罗 14:17}}；主也宣布：\Quote{人生活不只靠饼，也靠天主口中所发的一切言语。}此外，既然复活赐予我们与天使同等的生命，而天使没有食物，我们有充分理由相信，将像天使一样生活的人也会脱离这种功能。\par

% structure: p0136 source=h2 target=section reason=relative-heading-rank; base=h2

\section{XIX. \Emph{对那些说我们所盼的美好享受仍在于吃喝，因为经上记载人起初在乐园中靠这些生活的人。}}

1. 但或许有人会说，如果正如看来，我们从前靠吃而存在，而将来会脱离这一功能，那么人将不会回到同样的生活方式。然而，当我聆听圣经时，我不仅理解为身体的饮食或肉体的快乐；我也认出另一种食物，与身体的食物有某种相似，其享受只及于灵魂：\Quote{请你来吃我的食粮\BibleRef{箴 9:5}}，这是智慧对饥饿者的命令；主也宣布饥渴于这种食物的人是有福的，并说：\Quote{谁若渴，到我这里来喝}；伟大的依撒意亚也嘱咐那些能聆听他崇高之言的人\Quote{你们要喜乐地喝}。对于该受惩罚的人也有先知的警告，他们将遭受饥荒；但那\Quote{饥荒}不是缺乏饼和水，而是话语的断绝：\Quote{不是饥荒，不是口渴，而是听不到上主的话的饥荒\BibleRef{亚 8:11}}。\par

2. 那么，我们应当认为\Place{伊甸}的果子是配得上天主栽种的（\Place{伊甸}被解释为\Quote{欢乐}），并不怀疑人由此得到滋养；我们也不该认为乐园的生活方式是这种短暂易朽的食物：\Quote{园中各种树上的果子，}祂说，\Quote{你都可以随意吃\BibleRef{创 2:16}。}\par

3. 谁能给那有健康饥饿的人那棵在乐园中的树，它包含了一切美善，被称为\Quote{各种树}，这段经文赋予人分享它的权利？因为在这普遍而超越的宣告中，每一种美善的形式都与自身和谐，而整体是一。谁又能阻止我品尝那混杂而可疑的树上的果子呢？因为凡是略有眼力的人都清楚，那\Quote{各种}树是指什么，它的果子是生命；而那混杂的树，其结局又是死亡。因为那毫不吝啬地赐予对\Quote{各种}树的享受的主，必定是出于某种理由和预见，阻止人参与那些可疑的树。\par

4. 在我看来，我可以请伟大的\Person{达味}和智慧的\Person{撒罗满}作我这节经文的导师：因为两人都明白所允许的欢乐的恩宠是同一的——即那实际的美善，它实在是\Quote{一切}美善——\Person{达味}说：\Quote{你要在主内欢欣；}而\Person{撒罗满}则命名智慧本身（即主）为\Quote{生命树\BibleRef{箴 3:18}}。\par

5. 因此，这段经文赐食物给那按天主的肖像受造者的\Quote{各种}树，就是生命树；与这树相对的是另一棵树，它所给的食物是认识善恶的知识：——这并不是说它依次结出这些相反涵义的果子，而是它结出一种混杂着相反品质的果子，生命之主禁止吃它，而蛇却劝人吃，好为死亡预备入口；蛇的劝告得到了信任，它用美丽的外表和快乐的假象遮盖果子，使它悦目并刺激人想尝的欲望。\par

% structure: p0142 source=h2 target=section reason=relative-heading-rank; base=h2

\section{二十、 \Emph{乐园的生活是什么，以及什么是禁树？}}

1. 那么，那包含混杂在一起的善恶知识、并装饰以感官快乐的是什么呢？我想，我用\Quote{可知的}这个意义作为我思考的起点，并没有偏离目标。我认为，圣经在此用\Quote{知识}所指的并非\Quote{科学}；但根据圣经的用法，我发现\Quote{知识}和\Quote{辨识}之间有一定的区别：因为宗徒说，熟练地\Quote{辨识}善恶，是更完善状态和\Quote{锻炼过的感官}的标志，因此他也吩咐我们\Quote{要考验一切\BibleRef{得前 5:21}，}并说\Quote{辨识}属于属神的人\BibleRef{格前 2:15}；但\Quote{知识}并不总是理解为对某事的技能和熟悉，而是指对合乎心意之物的倾向——正如\Quote{主认识那些属于祂的人\BibleRef{弟后 2:19}}；祂对\Person{梅瑟}说：\Quote{我认识你超过众人；}而对于那些因邪恶而被定罪的人，那认识万有的说：\Quote{我从来不认识你们\BibleRef{玛 7:23}。}\par

2. 那么，结出这种混杂知识的果子的树，属于被禁止的东西；那果子由相反品质结合而成，有蛇来推荐它，这或许是因为：邪恶并不赤裸裸地暴露自己，它不以自己本来的面目出现——因为若邪恶不装饰某种美丽的颜色来诱使它所欺骗的人产生欲望，它必然不会得逞——但是，如今邪恶的本性在某种意义上是混杂的，将毁灭像罗网一样隐藏在深处，而在其欺骗性的外表下显示出某种好的幻影。财物的美丽在那些爱钱的人看来似乎是好的：然而\Quote{爱钱是万恶的根源\BibleRef{弟前 6:10}}；谁又会陷入放荡的恶臭泥沼中，若不是被这诱饵匆忙卷入激情的人认为快乐是美丽而可接受的？同样，其他的罪恶也隐藏它们的毁灭，初看似乎可接受，而某种欺骗使它们在轻率之人眼中被热切追求，代替了美善。\par

3. 既然多数人认为美善在于满足感官的事物，而实际的美善与看起来是\Quote{善}的有某种名称上的同一——因此，那朝向邪恶如同朝向美善的欲望，被圣经称为\Quote{善恶的知识；}\Quote{知识}，如我们所说，表达一种混合的倾向。圣经论到禁树的果子，不说它是绝对邪恶的（因为它装饰着美善），也不说它是纯粹美善的（因为邪恶潜伏其中），而是说它由两者组合而成，并宣称品尝它给接触它的人带来死亡；几乎大声宣告这教义：实际的美善在本性上是单纯而单一的，远离一切双重性或其反面结合，而邪恶则多彩多姿、外表华丽，被认为是一回事，而经验揭示为另一回事，对它的知识（即通过经验接受它）是死亡和毁灭的开始和先导。\par

4. 蛇看到了这一点，就指出了罪恶的恶果，不是明显地显示恶的本质（因为人不会被明显的恶所欺骗），而是给女人所见之物披上某种美的魅力，并在其滋味中注入感官享受的魔力，她似乎被他说服了：经上说：\BibleQuote{并且女人看见了}，经上说，\BibleQuote{那棵树上的果实实在好吃，又悦人眼目，且是可爱的；她就摘了果子来吃了}，那吃食成了人类死亡之母。这就是那种混合性质的果实，经文清楚地表达了那棵树被称为\Quote{能知善恶}的意义，因为就像掺了蜜的毒药，从它给人甜美的感官印象来看，它似乎是善的；但从它毁灭接触它的人来看，它是最恶的。因此，当邪恶的毒药对人生命发作了，人——那高贵者与名字，天主本性的肖像——就如先知所说的，成了\Quote{虚幻}。\par

5. 因此，肖像恰当地属于我们属性的较好部分；但我们生命中所有痛苦和可悲之事都远离对天主的肖似。\par

% structure: p0148 source=h2 target=section reason=relative-heading-rank; base=h2

\section{XXI. \Emph{论复活之所以被期待，与其说来自圣经的宣告，不如说来自事物本身的必然性}。}

1. 然而，邪恶并不能强大到胜过善的力量；我们本性的愚昧也不比天主的智慧更有力、更持久。因为那总是变动不居的，不可能比那永远不变且坚定于善的更坚固、更持久。但毫无疑问，天主的旨意具有不变性，而我们本性的可变性即使在邪恶中也不停留。\par

2. 那永远在运动的东西，如果它的进展朝向善，将永远不会停止向着面前的善前进，因为要走的道路是无限的——它不会找到所追求对象的任何界限，以致当它把握了对象后就最终停止运动。但如果它的偏向是相反的方向，那么当它完成了邪恶的历程并达到邪恶的极端时，那永远运动的东西，由于在邪恶中跑尽了可跑的长度后找不到其本性冲动的停留点，就必然将其运动转向善。因为邪恶并不延伸到无限，而是被必要的界限所包围，所以善似乎在邪恶的界限之后接连出现；因此，正如我们所说，我们本性的永远运动的特性最终会再次跑回善，从过去的灾祸记忆中学会谨慎的功课，以便它不再陷入同样的情境。\par

3. 因此，由于邪恶的本性被必要的界限所限制，我们的历程将再次存在于善之中。因为正如精通天文学的人告诉我们，整个宇宙充满了光，黑暗是由地球所形成的物体遮挡造成的；这种黑暗以圆锥体的形状被太阳光线隔离，根据球状物体的形状，在它后面。而太阳比地球大许多倍，从四面八方用它的光线包裹地球，在圆锥的极限处汇聚了光流；因此，如果（假设）有人能超过阴影延伸的尺度，他肯定会发现自己身在未被黑暗破碎的光中——同样，我认为我们应当这样理解我们自己：在越过邪恶的界限后，我们将再次生活在光中，因为善的本性与邪恶的尺度相比，是无限地超出的。\par

4. 因此，乐园将恢复，那棵树将恢复，那真正是生命树的树——肖像的恩宠和统治的尊威也将恢复。在我看来，我们的希望不是针对那些现在由天主交给人类用于生活必需之物，而是另一个国度，属于说不尽的奥迹。\par

% structure: p0153 source=h2 target=section reason=relative-heading-rank; base=h2

\section{XXII. \Emph{针对那些说\Quote{如果复活是卓越而美好的事，为什么它还没有发生，而是在某些时期被期待？}的人。}}

1. 然而，让我们关注讨论的下一点。也许有人，让他的思想展翅飞向我们盼望的甜蜜，认为我们没有更快地置身于那超越人的感官和知识的善境是一种负担和损失，对介于他和他的渴望对象之间的时间的延长感到不满。让他停止像一个孩子那样，因为延迟得到喜欢的快乐而烦恼；因为既然万物都由理性和智慧管理，我们决不能假设任何发生的事情没有其理性和其中的智慧。\par

2. 那么你会说，这理性是什么？按照它，我们痛苦生活的改变不会立刻发生，而是我们这沉重而有形体的存在等待某个确定的时间，直到万物的终结，那时人的生命仿佛从缰绳中释放，再次恢复自由，无拘无束地回到福乐和不受痛苦的境界？\par

3. 嗯，我们的答案是否接近事实的真相，\Quote{真理}本身可能清楚知道；但无论如何，我们理智所领会的是如下内容。那么，在我论证中，我再次拿起我们的第一个文本：天主说：\BibleQuote{让我们照我们的肖像，按我们的模样造人，天主于是造了人，照天主的肖像造了他。}\BibleRef{创 1:26} 因此，我们在普遍人性中所看见的天主的肖像，在那时就已完成；但那时亚当尚未存在；因为从尘土所形成的那东西被称为亚当，这是语源学的命名，正如那些熟悉希伯来语的人告诉我们的——因此那位在自己母语、以色列语言上特别有学问的宗徒，称那人为\BibleQuote{属土的}\BibleRef{格前 15:47}，\OriginalTerm{属土的}{\GreekText{χοϊκός}}，仿佛将亚当的名字翻译成希腊词。\par

4. 于是，人是按天主的肖像受造的；也就是说，那普遍的性体，那相似天主的物；不是整体的一部分，而是全部性体的圆满一起由全能的智慧如此造成。祂看见了——那将一切极限握在手中的那位，正如圣经告诉我们的，说\Quote{在祂手中是大地的所有角落}，祂看见了——那\Quote{知晓一切}甚至\Quote{在它们发生之前}那一位，以祂的知识涵盖它们，知道人性在其个体总和中将有多大的数目。但是，因为祂在我们受造的性体中察觉到了偏向邪恶的倾向，以及从与天使同等的地位自愿堕落后将与较低等的性体有份的事实，祂为此原因，在自己肖像中掺入了非理性因素（因为男女的区别并不存在于神圣和蒙福的性体中）；——我说，祂将非理性受造的独特属性转移到人身上，祂赐予我们种族增长的能力并非依照我们受造的崇高特性；因为祂并非在造那按自己肖像者时赐予人增长和繁衍的能力；而是当祂以性别区分它时，才说：\BibleQuote{你们要增长和繁衍，充满大地}\BibleRef{创 1:28}。因为这不属于神圣的，而属于非理性元素，正如历史所指示的，它叙述了这些话最初是由天主在非理性受造物的情况下说的；因为我们可以确信，如果在将男女区别印在我们性体之前，祂就已经赐予人通过这句话传达的增长能力，我们本不需要这种畜生繁衍的方式。\par

5. 如今，既然那些通过预知的活动预先被构想的人的总数将凭借这种动物性的生殖进入生命，天主——祂以一定的秩序和序列管理万有——由于我们性体向低于它之物的倾向（那同时看见未来和现在的一位，在它存在之前就已看见）使得某种这样的生殖方式对人类成为绝对必要——因此也预知了与人受造共延的时间，以便时间的长度能适应预定灵魂的进入，并且时间的流动和运动应在人性不再经由它产生的时刻停止；而当人的生殖完成时，时间应与它一同停止，然后万有的恢复会发生，并且随着世界的更新，人性也将从可朽坏和属地的变为不动情和永恒的。\par

6. 在我看来，神圣的宗徒考虑到了这一点，当他在写给格林多人的书信中宣布时间的突然停止，以及现在运动着的事物向相反终点的转变时，他说：\BibleQuote{看，我告诉你们一个奥秘；我们不全都要长眠，却全都要改变，在一刹那，在一眨眼之间，在最后的号角声中}\BibleRef{格前 15:51}。因为，如我所推测，当人性的全部补充达到了预定尺度的极限，因为不再有任何需要补充的灵魂数量，他教导我们，现存事物的改变将在时间的一个瞬间发生，将那没有部分或延伸的时间界限称为\Quote{一刹那}和\Quote{一眨眼}；这样，一个到达时间边缘的人（这是最后和最极端的点，因为没有什么欠缺以达到其极限）将不再可能通过死亡获得这种在固定时期发生的改变，而只会在复活号角吹响时，那号角唤醒死者，并将那些留在生命中的人，按照那些经历了复活改变的人的样式，立刻转变为不朽；这样，肉身的重量不再沉重，它的负担也不再将他们拖向大地，而是他们升空翱翔——因为他对我们说：\BibleQuote{我们将被提到云中，在空中与主相遇；这样我们就将永远与主在一起}\BibleRef{得前 4:17}。\par

7. 因此，让他等候那必然与人类发展同步的时期。因为就连亚巴郎和圣祖们，虽然渴望看到所应许的美物，不断寻求天上的家乡——如宗徒所说——至今仍处于盼望那恩宠的状态，因为保禄的话说：\Quote{天主为我们预备了更美的事，} \Quote{叫他们若不与我们同得，就不能完全。} \BibleRef{希 11:40} 如果他们仅凭信德和望德看见那些美物\Quote{从远处}、\Quote{且拥抱了它们} \BibleRef{希 11:13}，正如宗徒所见证的，他们之所以确信能享受所盼望的事，是因为\Quote{他们以为那应许者是可信的} \BibleRef{希 11:11}，那么，我们大多数人——或许因生活作风并未紧紧抓住那更美的盼望——又该怎样做呢？连先知的心灵也因渴望而憔悴，他在圣咏中承认这种炽爱，说他的\Quote{灵魂渴慕上主的庭院，}即使他必须被弃置于最卑微的位置，因为在那里居末位也比在此世不敬神的帐幕中居首位更伟大、更可羡慕；但他仍忍耐等候，诚然认为那里的生命是有福的，并视短暂参与它比\Quote{千日}更为可贵——因为他说：\Quote{在你的庭院一日，胜过千日}——然而他并不抱怨对现存事物的必要安排，认为人即使仅凭望德拥有那些美物已是足够的福乐；因此他在圣咏结束时说：\Quote{万军的上主，那信赖祢的人，真是有福！}\par

8. 因此，我们也不该因所盼望的迟延而烦恼，而应努力不要被排斥在盼望的对象之外。正如有人预先告诉一个没有经验的人说：\Quote{收割庄稼将在夏季来临，粮仓将装满，丰饶时期餐桌上将供应充足的食物，}那么，只有愚人才会急于催促果实的到来，而他本该播种并殷勤准备自己的庄稼；因为果实时期必定按指定的时间来到，不管他愿意与否；对于那些预先为自己获取了丰盛庄稼的人，和那些被果实时期发现毫无准备的人，它将显现得不同。我认为同样地，既然这宣告已清楚向所有人发出：改变的时期必将到来，那么我们的职责不是为时间而烦恼（因为祂说：\Quote{时候和日期不是你们应当知道的} \BibleRef{宗 1:7}），也不是追求那些必定削弱灵魂中复活盼望的计算；而是要将对所期望之事的信赖当作支柱，并藉着善行购买那将来的恩宠。\par

% structure: p0162 source=h2 target=section reason=relative-heading-rank; base=h2

\section{\Emph{二十三、承认世界存在开始的人，必然也同意其终结。}}

但如果有人观看世界当前的进程——其间的时间间隔按某种秩序进行——并说所预言的那些运动之停止不可能发生，那么此人显然也不相信起初天地是由天主创造的；因为承认运动有开始的人，无疑也不怀疑它也有终结；而不承认其终结的人，也不承认其开始；但正如我们藉着相信\Quote{我们因着信德，知道诸世界是因天主的话而形成的，}如宗徒所说，\Quote{以致看得见的是由看不见的所形成的，}同样，我们也必须用同样的信德来对待天主的话，当祂预言论及现存事物必然停止时。\par

2. 然而，关于\Quote{如何}的问题必须排除在我们的干预之外；因为即使在前述事例中，我们也是\Quote{藉着信德}承认看得见的是由尚未显现的事物形成的，而不去探究超出我们能力的事。尽管我们的理性在许多方面提出难题，对我们所信的事提供了不小的疑惑之机。\par

3. 因为在那事例中，好辩的人也可能藉着似是而非的推理推翻我们的信德，以致我们不认为圣经关于物质受造所传递的陈述是真实的——即一切现存之物从天主那里获得存在的开始。因为坚持相反观点的人主张物质与天主同永，并用自己的学说提出一些论据，例如：如果天主在本性上是单纯、非物质的，没有数量、大小或组合，并且排除了以形状来界定的概念；而一切物质则通过间隔来理解的广延来领会，并不能逃脱我们感官的知觉，而是在颜色、形状、体积、大小、阻力及其他附属属性中为我们所知——这些属性中没有一样能够设想在天主性内——那么，从非物质产生物质，从不伸展的\Quote{本性}产生有维度的\Quote{本性}，有什么方法呢？因为如果相信这些东西从那个源头获得存在，它们显然是在某种奥秘的方式中存在于祂之后才出现的；但如果物质的存在在祂内，祂如何能是非物质的而同时包含物质在祂内？同样地，所有其他标志物质\Quote{本性}的特征：如果数量存在于天主内，天主如何是没有数量的？如果组合的\Quote{本性}存在于祂内，祂如何是单纯的、无部分、无组合的？因此，论证迫使我们要么认为祂是物质的，因为物质从祂作为源头获得存在；要么，如果要避免这一点，就必须假定物质是由祂从外部\OriginalTerm{外来的}{ab extra}引入以创造宇宙的。\par

4. 如果这基础是外在于天主的，那么除了天主之外，必然还有某种与无生之天主同在永世中并存的东西；因此，这个论证假设了两个永恒且无生的存在，它们彼此共存——一边是作为工匠的天主，另一边是承受这种巧妙加工之物；如果有人因这论证的压力，为万物的创造者假定了一个物质基体，那么摩尼教徒会从他的特殊教义中找到何等支持，他以无生性将物质存在对立于善的存在。然而我们确实相信万物出于天主，正如我们听到圣经如此说；至于万物如何存在于天主内，这超越我们理性的问题，我们不去探究，因为我们相信万物都在天主能力范围内——既能使不存在的存在，也能按祂的旨意在其中注入特性。\par

5. 因此，既然我们认为天主圣意的能力足以成为万物从无中产生的因，那么我们也将以同样的能力归属于世界改造，而不会相信任何超出可能性的东西。此外，或许可以通过某种措辞技巧，说服那些就物质问题提出轻率异议的人，不要以为他们能对我们的陈述进行无可反驳的攻击。\par

% structure: p0168 source=h2 target=section reason=relative-heading-rank; base=h2

\section{第二十四章。 \Emph{驳斥那些说物质与天主同永的观点。}}

1. 因为归根结底，那种主张物质从可理解的、无形体的存在者获得其存在的观点，并非不符合常理。我们发现一切物质都由某些性质构成，若除去这些性质，物质本身就无法被观念把握。此外，在观念中，每一类性质都与基体分离；而观念是一种理智的而非形体的考察方法。例如，若有某种动物、树木或其他有物质存在的东西呈现在我们眼前，我们在理智讨论中会注意到许多关于基体的事物，其中每个观念都与我们观察的对象明确区分：颜色的观念是一回事，重量的观念是另一回事；同样，数量和某种独特的触觉性质的观念也是如此：因为\Quote{柔软}、\Quote{二肘长}以及我们提到的其他属性，在观念上既不彼此相连，也不与身体相连：每一个都根据其存在有其各自的解释性定义，与在基体中被观察到任何其他性质毫无共同之处。\par

2. 因此，如果颜色是可理解的，抵抗力也是可理解的，数量及其他类似性质也一样，而如果这些性质中的每一种都从基体中被抽离，那么身体的整个观念就解除了；那么似乎可以推论，我们可以假设这些事物的共现——我们发现它们缺位是身体解体的原因——产生了物质本性：因为正如一个没有颜色、形状、抵抗力、延伸、重量及其他性质的东西不是物体，而每种这些性质在其自身存在中却被发现不是物体而是物体之外的东西一样，反之，当指定的属性共现时，它们就产生了物体的存在。然而，如果对这些性质的感知是理智的事，而神性在本质上也是理智的，那么假设这些产生身体的理智诱因从无形体本性获得存在，就没有什么不协调：理智本性一方面赋予理智潜能以存在，而这些潜能的相互共现则产生物质本性。\par

3. 不过，让这段讨论作为题外话吧：我们必须再次将论述转向我们所接受的信仰，即宇宙从无中产生，并且当我们受圣经教导时，我们毫不怀疑它将来会转变为另一种状态。\par

% structure: p0172 source=h2 target=section reason=relative-heading-rank; base=h2

\section{第二十五章。 \Emph{即使是教外之人如何可以被引导相信圣经关于复活的教导。}}

1. 也许有人着眼于身体的分解，并以自己能力的大小来衡量天主，断言复活的想法是不可能的，说现在运动的事物不能静止，而静止的事物不能复活。\par

2. 然而，让这样的人以宣告复活者的可信性作为复活真理首要且最大的证据。因为所说的话的确信来自其他预言的结果：既然圣经发表了众多且不同的陈述，我们可以通过考察其他言论在虚假与真实方面的状况，来审视关于复活的教义。因为如果在其他事情上，陈述被发现是虚假的且没有实现，那么复活也不在真理范围之内；但如果所有其他陈述都有经验为其真理作证，那么基于它们，认为关于复活的预言也是真实的，似乎是合乎逻辑的。因此，让我们回忆一两个已经作出的预言，并将结果与预言相比较，以便我们能够借此知道这一想法是否具有真实的一面。\par

3. 谁不知道以色列民族昔日如何昌盛，兴起对抗世上一切权势？耶路撒冷城内的宫殿、城墙、塔楼、圣殿的宏伟建筑，又是什么样子？这些事物连主的门徒也似乎认为值得赞叹，以致他们怀着惊叹之情请主留意，正如福音史所载，说道：\Quote{这是何等伟大的工程，何等壮丽的建筑！}但主却向那些惊叹其现状的人预示那地方将来的荒凉和这美丽的消逝，说再过不多时，他们所见的必没有一块石头留在石头上。再者，在祂受难时，妇女们跟随祂，哀哭祂所受的不公判决——因为她们还不能看透所行之事的天主安排——但祂却要她们为祂所遭遇的事保持沉默，因为那不需要她们的眼泪，而要她们将哀号和痛哭留到真正该哭的时候，就是当那城被围困者包围，她们的苦难达到如此绝境，以致她们认为没有出生的人是有福的：在此祂也预言了那吞吃自己孩子的可怕行径，因为祂说在那些日子，子宫从不怀孕的才是有福的。\BibleRef{路 23:27} 那么，那些宫殿在那里？圣殿在那里？城墙在那里？塔楼的防御工事在那里？以色列人的权势在那里？他们不是几乎分散到全世界各地了吗？随着他们的覆灭，宫殿也一同被毁。\par

4. 在我看来，主预言这些和类似的事，并非为这些事本身——因为对听者而言，预言将要发生的事有何大益处？即使他们先前不知道将要发生的事，也会凭经验知道——而是为了藉此使他们对更重要的事产生信心：因为前者的实际见证也是后者真理的证明。\par

5. 正如一个农夫在解释种子的生命力时，若有一个不懂农事的人不信他的话，农夫只需从一斗种子中拿出一粒，展示其生命力，作为其余种子的凭据，便足以证明他的陈述——因为看见一粒小麦、大麦或斗中任何一粒种子，种在地里后长成穗子的人，便不会再对其他的种子存疑——同样，其他陈述所公认的确实性，在我看来，也足以作为复活奥秘的证据。\par

6. 然而，我们藉实际经验所认识的复活奇迹，更是如此；这经验主要不是靠言语，而是靠事实：因为复活的奇事既伟大又难以置信，祂便从较低级的奇迹能力开始，逐步训练我们的信心，使之能接受更伟大的事。\par

7. 正如一位母亲，起初以适当的关怀哺育婴孩，用乳房将奶水送到他柔嫩的口中；当他渐渐长大、长出牙齿时，她便给他面包，但不是坚硬难嚼的那种，免得粗糙的食物磨损柔嫩而未经锻炼的牙龈；她用自己的牙齿将面包软化，使之适合并易于食者的能力吸收；然后，随着他因成长而能力增强，她便逐步引导这习惯了柔软食物的婴孩，过渡到更坚固的营养；同样，主以奇迹滋养和培育人类心灵的软弱，有如一个尚未长成的婴孩，首先在一种绝望的疾病中，预演了复活的能力；这预演虽在成就上伟大，但其叙述却不至于令人不信：因为祂\Quote{斥责那热症}，就是烧灼西满岳母的凶猛热症，如此彻底地消除了病患，使那本已被认为濒临死亡的人，竟能\Quote{伺候}在场的人。\par

8. 接着，祂稍为增强能力。当贵族之子身处公认的死亡危险时（因为史载他快死了，其父喊道：\Quote{趁我的孩子还没有死，请你下来吧！}），祂再次使一个被认为快要死的人复活；这次，祂以更大的权能完成奇迹，甚至没有亲临现场，而是从远处以命令的力量赋予生命。\par

9. 再往后，祂上升到更高的奇事。因为祂在前往会堂长女儿家的路上，故意中途停留，公开了那患血漏妇女的暗中痊愈，好让死亡在这段时间内临到病人身上。当灵魂刚与身体分离，为这哀伤而恸哭的人正发出悲鸣时，祂以命令的话语使那少女复生，仿佛从睡中唤醒，以某种路径和次序，将人类的软弱引向更伟大的事。\par

10. 除此以外，祂以更令人惊奇的举动超越这些事迹，并以更崇高的能力举动为人们预备了相信复活的途径。圣经告诉我们，在犹太有个名叫纳因的城：那里有一个寡妇，她有一个独子，不再是孩童意义上的孩子，而是已经过了童年，进入成年阶段；叙述称他为\Quote{青年}。叙述用很少的话传达了很多：这叙述本身就是一场真正的哀悼；它说，死者的母亲\Quote{是个寡妇}。你看她不幸的分量，文本如何简短地展示了她受苦的悲剧？这句话是什么意思？意思是她没有希望再生儿子来弥补她刚刚在离去者身上遭受的损失；因为那女人是个寡妇，她没有能力指望另一个人来代替那个已去的人；因为他是她的独子；这表达了多么巨大的悲伤，任何不完全陌生于自然情感的人都能轻易看出。她只在痛苦中认识他，只用自己的乳房哺育他；只有他使她的餐桌充满欢乐，只有他是她家庭中光明的原因，在游玩、工作、学习、欢乐、游行、运动、青年聚会中；在母亲眼中，只有他是一切甜美和珍贵的事物。现在到了结婚的年龄，他是她种族的根基，是延续的嫩芽，是她晚年的支柱。此外，他生命时期的额外细节是另一场哀悼：因为称他为\Quote{青年}的人，是在述说他凋谢的美丽之花，说他刚刚用绒毛覆盖脸庞，还没有浓密的胡须，但脸颊的美丽依然明亮。那么，你认为，他母亲为他有多么悲伤？她的心会如何像被火焰吞噬一样；她会多么痛苦地延长对他的哀悼，拥抱着躺在面前的遗体，尽可能地延长对她的服丧，以便不加快死者的葬礼，而是让自己充分悲伤！叙述没有忽略这一点：因为耶稣\Quote{看见她}，经上说，\Quote{就动了怜悯的心}；\Quote{于是上前按住棺架，抬棺的人就停住了；}祂对死者说，\Quote{青年人，我对你说，起来 \BibleRef{路 7:13}}，\Quote{就把他活活地交给了他的母亲}。请注意，自从死者进入那种状态以来，并没有经过很短的时间，他几乎要被放入坟墓；主所行的奇迹更为伟大，尽管命令是相同的。\par

11. 祂的奇迹般的能力进展到更崇高的行动，使它的展示更接近人们怀疑的复活奇迹。主的一个同伴和朋友病了（病人名叫\Person{拉匝禄}）；主阻止了任何探望祂朋友的行动，尽管远离病人，以便在生命缺席的情况下，死亡可以借疾病找到空间和力量来做他自己的工作。主在\Place{加里肋亚}告诉祂的门徒\Person{拉匝禄}所遭遇的事，以及祂自己出发去他那里，要在他倒下时将他复活。然而，他们非常害怕，因为犹太人的愤怒，认为在那些想要杀害祂的人中间再次转向\Place{犹太}是一件困难和危险的事：因此，他们徘徊拖延，从\Place{加里肋亚}缓缓返回：但他们确实返回了，因为祂的命令占了上风，门徒被主带领，在\Place{伯达尼}被引入普遍的复活的前期奥秘。事情已经过去了四天；所有应尽的仪式都已为逝者举行；遗体被藏在坟墓中：它可能已经肿胀并开始腐烂解体，因为遗体在潮湿的泥土中腐烂，必然腐朽：这是一件令人回避的事，因为腐烂的遗体在自然约束下变得令人厌恶。在这一点上，被怀疑的普遍复活的事实通过一个更明显的奇迹得到了证实；因为没有人从重病中复活，也没有人从最后一息被救回生命——也没有刚刚死去的孩子复活，也没有即将被送往坟墓的青年从棺架上被释放；而是一个过了壮年的人，一具遗体，腐烂、肿胀，甚至已经处于解体状态，以至于即使是他自己的亲属也无法忍受主靠近坟墓，因为里面腐烂的遗体令人厌恶，通过一声呼唤就活了过来，证实了复活的宣告，也就是说，我们通过一个特定的经验学习到对普遍复活的期待。因为正如在宇宙的更新中，宗徒告诉我们：\Quote{主亲自要从天降下，有呼喊声，有总领天使的声音 \BibleRef{得前 4:16}}，并用号声使死人复活，成为不朽的——同样，现在那在坟墓里的，听到命令的声音，就摆脱死亡如同摆脱睡眠，并摆脱了临到遗体状态的腐朽，完整健康地从坟墓中跃出，甚至没有被缠在手脚的殓布束缚所妨碍。\par

十二、这些事难道太小，不足以在死人复活的事上产生信德吗？或者你寻求你的判断在这点上应由其他证据来确认？实际上，在我看来，主没有白白地对葛法翁的人说话，当时祂仿佛以人的身份对自己说：\Quote{你们必定要对我说这俗语：\InnerQuote{医生，医治你自己吧！}}因为祂在习惯使人在其他身体上看到复活的奇迹之后，理应在自己的人性中证实祂的话。你看见那所宣告的事在别人身上运作——那些即将死去的人，刚停止呼吸的孩童，坟墓边缘的青年，腐烂的尸体，都因一个命令而同样恢复生命。你寻找那些因伤口和流血而死亡的人吗？难道生命赐予之力的任何软弱阻碍了他们身上的恩宠？请看那双手被钉子穿透的祂，请看那肋旁被长矛刺穿的祂；用你的手指探入钉痕，将你的手探入枪伤；你当然能猜出枪尖可能到达多深，如果你按照外部伤疤的宽度来推算向内进入的深度；因为允许人手探入的伤口，显示了铁器进入的深度。如果祂已经复活，我们大可以发出宗徒的惊叹：\Quote{\BibleQuote{怎么你们中有人说没有死人复活呢？}\BibleRef{格前 15:12}}\par

十三、既然主的每一预言都由事件的结果证明为真实，我们不仅从祂的话中得知了这一点，而且从那些借复活返回生命的人身上，在行动中领受了这预言的证明，那么还有何机会留给那些不信的人呢？我们岂不应该向那些借着\Quote{哲学和虚妄的欺骗\BibleRef{哥 2:8}}歪曲我们纯朴信德的人告别，并持守我们纯洁的宣认，借先知的话简要学习恩宠的方式？他的话是：\Quote{祢收回他们的气息，他们就归去，归于他们原来的尘土。祢发出祢的圣神，他们便被创造，祢使地面更新；}那时他也说，主因祂的作为而喜乐，罪人已从地上消灭：因为当罪本身不再存在时，谁能被称为罪之名呢？\par

% structure: p0186 source=h2 target=section reason=relative-heading-rank; base=h2

\section{第二十六章　复活并非不可能}

一、然而，有些人因人类推理的软弱，以我们的能力来衡量天主的德能，坚持认为超出我们能力的事，对天主也是不可能的。他们指出古时死者的消失，以及被火化为灰烬者的遗骸；此外，他们还想象肉食野兽，以及那吞下遇海难水手身体后又成为人食、并经过消化进入食者体内的鱼；他们列举了许多此类琐碎之事，不配于天主伟大的德能与权威，以此推翻教义，争论说天主不能通过同样的途径恢复人自己的东西。\par

二、但我们通过承认身体分解为其组成部分确实发生，并且不仅土归土，正如神圣之言所说，而且气和湿气也回归同类的元素，我们每个组成部分都回归其所属的本性，从而简要地截断了他们逻辑愚昧的长篇大论。尽管人体被分散于肉食鸟类或最野蛮的野兽，成为它们的食物，尽管它落入鱼齿之下，尽管它被火化为蒸气和尘埃，无论人在论证中假设人被移往何处，他确实仍留在世界中；而世界，正如灵感之言告诉我们，是在天主手中。那么，如果你对自己手中的任何东西都不无知，你会认为天主的认知比你自己的能力更软弱，以至于无法发现神圣范围之内最微小的事物吗？\par

% structure: p0189 source=h2 target=section reason=relative-heading-rank; base=h2

\section{第二十七章　当人体分解为宇宙元素时，各人能从共同源头恢复自己的身体是可能的}

一、然而你或许会考虑到宇宙的元素，认为当我们体内的气已回归其同类元素，热、湿气和土性也各自与自己的同类混合后，从共同源头回归属于个体的自身，是一件困难的事。\par

二、你难道不通过人的例子来判断，即使这一点也未超出天主德能的界限吗？你必定曾在某处见过从各处收集来的某种动物的共同畜群；但当它再次分归其主人时，对家的熟悉以及家畜身上的标记有助于恢复各人自己的。如果你也这样想象自己，你就离正道不远了：因为灵魂倾向于依附并渴慕与它结合的身体，由于它们的混合，某种隐秘的亲密关系和认知能力也附着于它，仿佛某些标记由本性印刻，借助这些标记，共同体保持不混淆，由区别性记号分开。现在，既然灵魂再次吸引那属于自己并适当属于它的东西，请问，对于天主的德能，有什么困难能阻碍同类事物的汇聚，当它们由本性不可言喻的吸引力被催向自己的位置时？因为即使在分解之后，我们复合本性的某些标记仍留在灵魂中，这在阴府中的对话\BibleRef{路 16:24}中显示，那里尸体已被送入坟墓，但某些身体标记仍留在灵魂中，借此拉匝禄被认出，富人也不被不知。\par

3. 因此，相信在復活的身體中，個人的成分會從共同庫存中返回，尤其是對任何仔細審視我們本性的人而言，並非沒有道理。因為我們的存在並非完全處於流變和變化之中——因為那在本性上毫無穩定的東西，當然完全無法理解——但按照更準確的說法，我們組成部分中的某一部分是靜止的，而其餘部分則經歷變化的過程：身體一方面通過生長和縮小而改變，像衣服一樣更換其不同身材的覆蓋物，而另一方面，形式本身在一切變化中保持不變，不偏離自然最初賦予它的標記，而是在身體所經歷的一切變化中，以其自身的同一性標誌顯現。\par

4. 然而，我們必須從此論述中排除因疾病而發生的形式變化：因為疾病的畸形像某種奇怪的面具佔據了形式，而當藉著話語將其除去時，如同在 \Person{纳阿曼}（叙利亚人）或福音書中所記載的那些人的情況，被疾病隱藏的形式便藉著健康再次以其自身的同一性標誌顯現出來。\par

5. 現在，在我們靈魂中相似於\OriginalTerm{天主}{God}的那個元素，與其說是通過變化而處於流變和變化中的東西，不如說是我們組成中這個穩定不變的元素與之相連：既然組合的各種差異產生了形式的多樣性（而組合無非是元素的混合——我們所說的元素是指為宇宙的製造提供基礎的那些元素，人的身體也由它們構成），而形式必然像印章的印記一樣留在靈魂中，那麼那些從印章接受了其印記印象的事物，就不會不被靈魂所辨認，而是在世界重整之時，靈魂重新接回一切與形式之印記相應的事物：確實，一切最初由形式所印記的事物都會與之相應；因此，個體所固有的東西從共同源頭再次返回它，並非沒有道理。\par

6. 據說水銀若從盛裝它的容器中傾倒在塵土斜坡上，會形成小珠並散落在地，不與所遇到的任何物體混合；但如果有人將四處分散的物質收集到一處，它就會流回與其同類的物質，只要沒有什麼障礙物阻止它與同類混合。我認為，關於人的複合本性，我們也應類似地理解：只要\OriginalTerm{天主}{God}賜予它這能力，各部分就會自發地與屬於它們的部分結合，而不會對重整它們本性的天主造成任何阻礙。\par

7. 此外，在從土地生長的植物中，我們沒有觀察到自然在小麥、小米或任何其他穀物或豆類的種子上花費任何勞動，將它們變成秸稈、穗或穀粒；因為適當的養分無需費力，就從共同源頭自發地流向每顆種子的個體性。那麼，如果供給所有植物的水分是共同的，而每一棵由它滋養的植物都為其生長吸取應有的供應，那麼在復活的道理中，如同種子的情況一樣，每位復活者吸引屬於自己的東西，又有什麼新奇之處呢？\par

8. 這樣，我們可以從各方面學到，復活的宣講並不包含任何超出我們通過經驗所知的事實。\par

9. 然而，我們尚未提及關於我們自身最顯著的一點：我指的是我們存在的開端。誰不曉得自然的奇跡——母親的子宮接受的是什麼——它產生的是什麼？你看，植入子宮作為身體形成之開端的東西，在某種意義上是單純而同質的；但有什麼語言能表達所形成的複合身體的多樣性呢？而且，若不是在自然中普遍學到這樣的事，誰會認為實際發生的事情是可能的——那微小而不起眼的東西竟是如此偉大之事的開端？我所說的偉大，不僅指身體的形成，更指比這更奇妙的事，即靈魂本身，以及我們在其中所見的種種屬性。\par

% structure: p0199 source=h2 target=section reason=relative-heading-rank; base=h2

\section{廿八、\Emph{反駁那些說靈魂先於身體存在，或身體先於靈魂形成的人；其中也駁斥關於\OriginalTerm{靈魂轉世}{transmigration of souls}的虛構故事}。}

1. 或許討論教會中提出的關於靈魂與身體的問題，並未超出我們當前的主題。我們之前的一些處理\Quote{原理}問題的人，認為應該說靈魂在一個自己的群體中作為一個民族先已存在，而且他們之中也有善惡的標準；那在善中持守的靈魂，不會經歷與身體的結合；但若它脫離了與善的共融，就會墮落到這個較低的生命中，因此進入身體。相反，另一些人依據梅瑟所記述的人受造的順序，說靈魂在時間次序上後於身體，因為天主先從塵土取土塑造了人，然後藉著祂的噓氣使這被塑造者有了生命；藉此論證他們證明肉體比靈魂更高貴——先前形成的比之後注入的更高貴：因為他們說靈魂是為身體而造的，使被塑造者不致沒有氣息和運動；而且，凡為他物而造的東西，肯定不如那為它所造的東西寶貴，正如福音告訴我們說：\BibleQuote{靈魂比食物貴重，身體比衣服貴重}\BibleRef{瑪 6:25}，因為後者是為了前者的緣故而存在——靈魂不是為食物造的，我們的身體也不是為衣服造的；而是當前者已經存在時，後者才為它們的需要而提供。\par

2. 既然这两种理论所涉及的教理——即一方面那些将灵魂归之于一种神话般先行存在的特殊状态的人，以及另一方面那些认为灵魂是在身体之后被创造的人——的教理备受批评，或许有必要对这两种教理所包含的陈述逐一加以审查；然而，若要从正反两方面彻底介入并驳斥这些理论，并揭示其中所包含的种种荒谬之处，则需要大量的论证和时间；不过，我们将尽可能简要地审视上述两种观点，然后继续我们的主题。\par

3. 坚持前一种理论、断言灵魂的状态先于其在肉身中的生命的人，在我看来似乎并未摆脱异教关于灵魂轮回的神话理论；因为若仔细审视，便会发现他们的理论必然导致这一点。他们告诉我们，他们的一位贤人曾说过，他作为同一个个体，曾生为人，后来变成女人的形状，与鸟类一同飞翔，长成一丛灌木，并获得了水生动物的生命——依我判断，说这些话的人并未远离真理：因为诸如说同一个灵魂经历了如此多的变化这类理论，确实只适用于青蛙或寒鸦的聒噪，或鱼类的愚钝，或树木的无感觉。\par

4. 这种荒谬的原因在于——灵魂先存的假设，因为这种理论的第一原则通过推论引向下一步和相邻阶段，直到它达到这一点而令人惊异。因为如果灵魂由于某种邪恶而与更崇高的状态分离，在尝过肉身生活之后（正如他们所说）再次成为人，并且如果肉身生活（正如可以推断的那样）相比永恒和非肉身的生活被认为更易受激情影响，那么自然的结果是：处于这样一种包含更多犯罪机会的生活中的灵魂，既被置于更大的邪恶区域，也比以往更易于受激情影响（而在人类灵魂中，激情是对非理性状态的相似性的顺从）；并且当它与之紧密联系时，它便下降到兽性；一旦它通过邪恶踏上这条路，即使在非理性状态下也不会停止向邪恶前进：因为停止邪恶是走向德行的冲动的开始，而在非理性之物中德行并不存在。因此，它必然不断变得更糟，总是走向更堕落的状态，总是发现比它目前所处的状态更坏的东西；正如感官本性低于理性本性一样，从感官本性下降到无感觉的本性也是如此。\par

5. 到目前为止，他们的理论——即便它越过了真理的界限——至少是通过一种逻辑顺序从一个荒谬推出另一个荒谬：但从这一点开始，他们的教导便呈现出不连贯的神话形式。严格的推论指向灵魂的完全毁灭；因为一旦从崇高状态堕落，灵魂将无法在邪恶的任何程度上停留，而是通过与激情的关系从理性转向非理性，并从后者转向植物的无感觉；而在无感觉之后紧接的，可以说，是无生命；在此之后随之而来的是不存在，因此通过这种推理，他们完全将灵魂推入虚无：这样，它就不可能再次回到更好的状态：然而他们却使灵魂从一丛灌木返回为人：因此他们证明了灌木中的生命比非肉体的状态更宝贵。\par

6. 已经表明，灵魂中发生的恶化过程可能会向下延伸；而在无感觉之下我们找到无生命，他们的理论原则通过推论将灵魂带到那里：但是因为他们不愿意这样，所以他们要么将灵魂排除在无感觉之外，要么，如果要使灵魂回到人类生活，就必须如所说那样，宣称树木的生命优于原始状态——如果堕落朝向邪恶是从前者发生，而回归德行是从后者发生的话。\par

7. 因此，他们这种认为灵魂在肉身生活之前有其独立生活、并因邪恶而被束缚于身体的理论，既没有起点也没有终点：至于那些断言灵魂是在身体之后被创造的人，他们的荒谬性已在上文得到证明。\par

8. 因此，两种理论同样应被拒绝；但我认为我们应当在这两种理论之间引导我们自己的教理走向真理：这个教理就是，我们不应按照异教的错误，认为灵魂随着宇宙的运动而旋转，由于某种邪恶而沉重，因无法跟上球体运动的迅疾而坠落于地。\par

% structure: p0208 source=h2 target=section reason=relative-heading-rank; base=h2

\section{二十九. \Emph{确立灵魂与身体存在的原因同一个的教理。}}

1. 我们的教义也不应当从将人塑造成泥像开始，并说灵魂是为了这个而存在的；因为果真如此的话，理智本性就将显得不如泥像宝贵了。但是，既然人是一个，是由灵魂和身体构成的存有，我们就应当假定他的存在的开端是单一的，为两部分所共有，这样他就不会在自身中发现先后之分——如果身体元素在时间上在先，而另一个随后添加的话；相反，我们要说，在天主预知的能力中（按照我们此前在论述中稍早提出的教义），人性的全部丰满已经有预存（对此先知著作作证，它说天主\Quote{知道一切，在它们存在之前}），而在个体的受造中，不应将一方置于另一方之前，既不是灵魂在身体之前，也不是相反，这样人就不致因时间上的差异而分裂，与自己为敌。\par

2. 因为按照宗徒的教导，我们的本性被视为双重的，由可见的人和隐藏的人构成，如果一方在先，另一方继之而来，那么创造我们的那位的德能将显明在某种方式上是不完美的，仿佛不能同时完全胜任全部工作，而是将工作分开，先后忙于每一半。\par

3. 但是，正如我们在麦子或其他任何谷粒中所说的，整个植物的形态都潜在地包含在内——叶子、茎、节、谷粒、芒——我们在描述其本性时，并不说其中任何一样有预存，或在其他之前产生，而是说存留于种子内的能力按照某种自然次序显现，绝非有另一种本性注入其中——同样，我们认为人的胚芽拥有其本性的潜能，从一开始存在就与它一同被播种，并通过一个自然序列展开和显现，逐渐达至其完美状态，不借助任何外在之物作为通往完美的阶梯，而是自身按着适当次序将自己提升至完美境界；因此，既不能说灵魂在身体之前存在，也不能说身体没有灵魂存在，而是说二者有一个共同的开端，这个开端按照天上的观点，是在天主原始意志中奠定的；按照另一种观点，则是在生育之际开始存在的。\par

4. 因为正如在植入以形成身体的东西中，在它开始成形之前，我们不能分辨肢体的关节，同样，我们也不能感知其中灵魂的诸特性，直到它们开始运作；并且正如无人会怀疑所植入的东西被塑造为不同的肢体和内脏器官，不是通过从外部引入任何其他能力，而是通过存在于其中的能力将自身转化为这种能量的显现——同样，我们也可通过类似的推理，在灵魂的情况中同样假定，即使它没有通过任何活动的显现而被视觉识别，它仍然在那里；因为未来之人的形态也是潜在地在那里，但由于在必要的顺序允许它被看见之前不可能使其显现而被隐藏；同样，灵魂也在那里，即使不可见，也将通过它自身特有的和本性的运作而显现，随着身体的生长而同步进展。\par

5. 因为既然受孕的潜能不是从死去的身体中分泌出来，而是从有生命和活着的身体中分泌出来，我们因此断言，我们不应认为从活体发出的、成为生命之缘由的东西本身是死的和无生命的；因为在肉体中，无生命的东西肯定是死的；而死亡的状态是由灵魂的离去而产生的。那么，在这种情况下，一个人岂不是在断言离去先于拥有——如果他坚持认为作为死亡状态的无生命状态先于灵魂的话？如果任何人寻求更清晰的证据，证明那成为受造生物形成之开端微粒的生命，也可以从其他迹象中获得关于此点的概念，通过这些迹象，有生命的东西与死亡的东西得以区分。因为对人而言，我们认为温暖、活动和运动是生命的证据，而在身体中，寒冷和静止状态无非就是死亡。\par

6. 既然我们看到我们所谈论的是温暖和活动的，我们由此进一步推断它并非无生命的；但是，正如就其身体部分而言，我们不说它是肉、骨头、头发以及我们在人身上观察到的一切，而是说它潜在地是这些东西中的每一个，却没有显眼地显现为这样；同样，对于属于灵魂的部分，理性、欲望、愤怒以及灵魂的一切能力的元素尚不可见；然而我们断言它们在其中有自己的位置，并且灵魂的诸能量也随着主体以类似于身体形成和完善的方式成长。\par

7. 因为正如一个人完全发育时，灵魂有明显的标记活动，同样，在他存在的开始，他在自身中显示出那种适合并符合他当下需要的灵魂的合作，即通过植入的物质为自己准备合适的居所；因为我们不认为灵魂能适应一个陌生的建筑，正如我们不可能将印在蜡上的印戳适配到与其不一致的雕刻上。\par

8. 因为身体从极小的起点发展到完美的状态，灵魂的活动也是如此，因应主体的成长而增长，与其一同壮大。最初的形成，首先只有生长和营养的能力，如同埋在地下的根；因为接受者的有限本性不允许更多；然后，当植物出现在光中，向太阳展示嫩芽时，感觉的恩赐也绽放出来，但最后，当它成熟并长到适当高度时，理性的能力开始像果实一样闪耀，不是一下全部展现出它的全部活力，而是通过照顾，随工具的完美而增长，总是结出与主体能力相称的果实。\par

9. 然而，如果你试图追溯灵魂在身体形成中的运作，\Quote{你要谨慎自己\BibleRef{申 4:23},} 如梅瑟所说，你会如同在书中读到灵魂工作的历史；因为本性本身比任何论述更清晰地为你解释灵魂在身体中的各种活动，无论是总体还是个别构造的行为。\par

10. 但我认为用言语详细说明在我们自身内可以发现的东西是多余的，仿佛我们在解释某些超越我们边界的奇迹：——谁看自己还需要言语教导他自己的本性？因为一个人思考自己生活的方式，并了解身体在每一种生命活动中有多么密切相关，就可以知道灵魂的植物性原理在最初形成那个开始存在的东西时做了什么；因此，这也使那些关注此事的人清楚，为了产生生命存在而从活体分离出来植入的东西，在本性的工坊中并非死物或无生命的。\par

11. 此外，我们把果核和从根上撕下的部分种在地里，它们并未因死亡而失去内在的生命力量，而是暗中却确实活着保存了原型的特性；围绕它们的土壤并非从外部注入这种力量（因为那样连死木也会生长），而是使它们内在的东西显现出来，用自身的湿气滋养它，使植物完美，成为根、树皮、髓和枝干，这不可能发生，除非一种自然的力量与它们一同植入，这种力量从周围环境中吸取同类的适当营养，成为灌木、树木、麦穗或某种灌木类植物。\par

% structure: p0220 source=h2 target=section reason=relative-heading-rank; base=h2

\section{XXX. \Emph{从医学观点略论我们身体的构造。}}

1. 我们身体的精确构造，每个人通过视、光和感知的经验自我教导，以自身的本性为导师；任何人也可以准确了解一切，只要查阅那些精通此类事务的人在书中所做的研究。在这些作者中，有些通过解剖了解我们个别器官的位置；另一些思考并解释了身体所有部分存在的原因；因此由此产生的人体知识对学生来说已经足够。但如果有人进一步寻求教会成为他在这一切点上的导师，以便他无需从外部声音那里得到任何东西（因为属灵的羊群的习惯，如主所说，就是不听从陌生人的声音），我们也将简要处理这些事。\par

2. 关于我们身体的三种事物，为了它们我们特定的部分被形成。生命是有些部分的原因，美好的生活是其他部分的原因，另一些部分则是为后代的延续而设计的。我们体内所有那些没有它们人类生命就不可能存在的事物，我们认为分为三部分：脑、心和肝。此外，所有那些额外祝福、本性的慷慨，借此她赐予人好好生活的恩赐，是感官器官；因为这类事物并不构成我们的生命，因为即使缺少其中一些，人往往仍处于生命状态；但没有这些活动形式，就不可能享受生活中的快乐。第三个目的是关于未来和生命的延续。除这些外，还有某些其他器官，它们共同帮助维持生命的延续，通过自身的方式输入适当的补给，如胃和肺，后者通过呼吸扇动心脏的火，前者为内脏引入营养。\par

3. 因此，我们的构造如此划分，我们必须仔细注意，我们的生命能力并非仅由单一器官以一种方式支持，而是本性在将生存的手段分配于几个部分时，使每个个体的贡献对整个整体成为必要；正如本性为了生命的安全和美丽而设计的事物也是多种多样，彼此差异很大。\par

4. 然而，我认为应当先简要讨论构成我们生命的事物最初的开始。至于整个身体的物质，它作为特定肢体的共同基体，目前可以不加评论；因为关于一般自然物质的讨论，对我们考察部分的目的不会有任何帮助。\par

5. 既然所有人都承认，我们体内有我们在宇宙中看到的所有元素——热与冷，以及另一对性质：湿与干——的一份，我们必须分别讨论它们。\par

6. 于是我们看见，控制生命的能力有三种：第一种借着其热力产生普遍的热度，第二种借着其湿度使被温暖的东西保持潮湿，使得活物处在一种中间状态，这是由于对立本性中每一种的特质所施加的力量之均衡（潮湿元素不会被过多的热力烤干，热力元素也不会被过多的湿气扑灭）；第三种能力则凭其自身作用，将各个部分以某种一致与和谐维系在一起，用其自身提供的纽带连接它们，并将那自我推动和决定的力量发送到它们里面，这力量若失效，那部分就变得松弛而失去活力，因为没有那决定性的精神而被遗弃。\par

7. 或者更确切地说，在处理这些问题之前，我们应当先留意自然在身体的实际构造上的巧妙工艺。因为坚硬和抗拒的东西不容感觉的作用（正如我们可见于我们自己的骨头，以及地上的植物，我们注意到其中确实有某种生命形式，因为它们生长并接受营养，然而其物质的抗拒特性不允许它们有感觉），为此原因，有必要提供某种像蜡一样的构造，可谓之为此，以供感觉作用，它具有被能触动感觉的事物之印记所印刻的能力，既不会因过多水分而变得模糊（因为印记不会留在潮湿物质上），也不会因异常的坚固而抵抗（因为不易屈服的东西不会从印象中得到任何印记），而是处于软硬之间的状态，好使活物不至于缺少自然最美妙的一切运作——我指的是感觉的运动。\par

8. 现在，如果柔软而易弯曲的物质没有来自硬质部分的辅助，那么它无疑会像软体动物一样，既没有运动也没有关节；因此，自然在身体中混合了骨头的硬度，并用紧密的连接将其相互结合，通过筋腱将其关节连接在一起，这样就将那能接受感觉的肌肉围裹在周围，使其表面比原本应有的更坚硬和更紧绷。\par

9. 于是，将整个身体的重量支撑在骨头的这种物质上，如同支撑在承载建筑重量的柱子上一样，她并未将骨头无分割地植入整个结构中；因为那样的话，人就会保持不动和没有活动，如果他这样被构造，就像一棵树固定在一处，既没有双腿的交替运动来前进其运动，也没有双手的服务来照料生活的便利；但现在我们看见，她通过这种装置，使工具能够行走和工作，她借着通过神经延伸的决定性精神，将冲动和运动的能力植入身体中。由此产生了双手的服务，如此多样而多变，与每一个思想相呼应。由此产生了，仿佛是某种机械装置一般，颈部的转动、头部的弯曲和抬起、下巴的动作、眼睑随思想而进行的开合，以及其他关节的运动，这都是通过某些神经的收紧或放松而进行的。贯穿这些的能力表现出一种独立的冲动，通过某种自然的管理，与意志的精神一起在每个特定部分工作；但一切的根本，以及神经运动的原理，存在于围绕大脑的神经组织中。\par

10. 我们认为，当运动的能力被显示在此处时，我们无需花更多时间探究这样一个东西存在于哪个生命器官中。但是，大脑以特殊方式有助于生命，这由相反情况的结果清楚显示：因为如果围绕它的组织受到任何伤口或损伤，死亡立即随之而来，自然连一刻也无法忍受这种伤害；正如地基被抽走时，整个建筑连同那部分一起倒塌；而那器官，其损伤显然导致整个活物的毁灭，可被恰当地承认为包含生命的原因。\par

11. 而且，在那些已经停止生活的人中，当植入我们本性的热力熄灭时，那已死的东西变冷，我们由此也承认生命的原因在于热力：因为我们必须必然承认，活物因那东西的存在而存续，那东西一旦缺失，死亡的状态就发生。我们理解心脏就是这种力量的源头和原理，因为它有如管道般的通道，互相生长并有许多分支，将温暖而热烈的精神扩散到全身。\par

12. 既然自然也必须为热力元素提供某种营养——因为火不可能独自持续，没有适当的食物来滋养它——因此，从肝脏如源头流出的血液通道，伴随着温暖的精神在整个身体的途中，使得彼此不会因孤立而成为疾病并破坏结构。让这教导那些超越公平界限的人，因为他们从自然中学到贪婪是一种导致毁灭的疾病。\par

13. 但是，由于唯有神性不需要任何东西，而人的贫乏需要外在的帮助以维持其存在，因此自然，除了那三种我们说过调节整个身体的能力之外，还从外部引入输入的物质，通过不同的入口引入适合于那些能力的东西。\par

14. 因为对于血液的源头，即肝脏，她通过食物来供应：因为不时以这种方式输入的物质，预备从肝脏涌出的血液之泉，如同山上的雪以其自身的湿气增加低处的水泉，将其液体强行渗透到下方的血管中。\par

15. 心脏中的呼吸由邻近的器官提供，这器官称为肺，是空气的容器，通过插入其中的气管（延伸至口腔）从外部吸入气息。心脏位于肺的中间（它本身也不停地运动，模仿永远运动之火的活动），像铁匠铺的风箱一样，从邻近的空气中吸引供应，通过扩张填充其凹陷，同时煽动自身的火元素，向邻近的管道吹气；它不停地进行这一过程，通过扩张将外部空气吸入自身凹陷，通过压缩将空气从自身注入管道。\par

16. 在我看来，这就是我们这种自主呼吸的原因；因为常常当心智忙于与他人交谈，或身体在睡眠中放松而完全静止时，气息的呼吸并未停止，尽管意志并未对此予以合作。现在我推测，由于心脏被肺包围，且其结构后部附着于肺，通过自身的扩张和压缩运动肺，吸气与呼气正是由肺完成的：因为肺是轻质多孔的身体，其所有凹陷都开口于气管基部，当它们收缩并受压时，必然通过压力挤出残留在腔内的空气；而当它们扩张并打开时，则通过伸展将空气吸入真空。\par

17. 这便是这种不随意呼吸的原因——火元素不可能保持静止：因为运动本是热的功能，而我们知道热的根源在于心脏，在此器官中持续进行的运动便产生了通过肺的不断吸气和呼气；因此，当火元素异常增强时，那些发热者的呼吸变得更加急促，仿佛心脏正试图通过更猛烈的呼吸来扑灭植入其中的火焰。\par

18. 但由于我们的本性贫乏，需要从各方面获得维护自身的供应，它不仅缺乏自身固有的空气和激发热量的气息——这些需要从外部输入以维持生命——而且用以填充身体比例的养分也是输入之物。因此，它通过饮食补充缺乏，在身体中植入某种能力，用以吸收所需之物并排斥多余之物；为此，心脏之火也给本性以不小的帮助。\par

19. 因为，根据我们之前的描述，心脏用其温暖的气息点燃各个部分，是最重要的生命器官。我们的造物主使它以其有效的能力在所有点位上运作，不让任何部分对整体机体的调节无效或无益。因此，在后方，它进入肺，并通过持续的运动将肺牵向自身，扩张通道以吸入空气，再压缩它们从而将被困空气呼出；而在前方，它附着于胃上端的空间，温暖它并使其以运动回应自身的活动，唤醒它，不是为了吸入空气，而是为了接收合适的食物：因为气息和食物的入口彼此靠近，纵向并排延伸，并在上端以同一界限终结，因此它们的口部相邻，通道在同一个口中结束，通过其中一个实现食物的进入，通过另一个实现气息的进入。\par

20. 然而在内部，通道连接的紧密性并非处处保持；因为心脏介入二者基部之间，向一个注入呼吸的力量，向另一个注入营养的力量。火元素自然倾向于寻求作为燃料的物质，营养的容器必然发生这种情况；因为火通过邻近的温暖越渗透它，它就越将滋养热量的东西吸引到自己身上。我们将这种冲动称为食欲。\par

21. 但是，如果容纳食物的器官获得了足够的材料，即使如此，火的活动也不会静止：而是像在铸造厂中一样，对材料进行某种熔化，溶解固体，将其倾倒并转移，宛如通过漏斗，到邻近的通道：然后将粗糙物质与纯净物质分离，将精细部分通过某些管道送至肝脏入口，并将食物的沉淀物排至肠的更宽通道，通过让其在其众多弯曲中翻转，将食物在肠中保留一段时间，以免如果通过一条直通道轻易排出，会立即再次引起动物的食欲，而人，如同无理性的动物种族，可能永不停歇于这类营生。\par

22. 然而，正如我们所看到的，肝脏特别需要热量的合作以将液体转化为血液，而此器官位置远离心脏（因为我想，若它本身是生命力的一个本源或根源，就不可能与另一个这样的本源邻近而不受阻碍），为了使系统不会因生热物质放置的距离而受损，一个肌肉通道（精通此类事物者称之为动脉）从心脏接收加热的空气并将其输送至肝脏，在其开口处开在靠近液体进入点的某处，并通过其热量温暖潮湿物质，将某种类似火的东西与液体混合，并使血液因所产生的火色而呈现红色。\par

23. 由此再发出，某些双通道，各自如管子般包围自己的水流，将空气与血液分散至全身各处（为使液体物质得享自由运行，当它伴随热质运动而变得轻盈时）；同时，生命力的两个原则相互混合（那分散热的，与那供应水分给全身各部的），它们仿佛从所处理的物质中，对生命经济中的最高力量作出某种强制性的奉献。\par

24. 如今，这力量就是那被认为存在于脑膜和脑中的力量；由此而来，关节的每一运动、肌肉的每一收缩、对各个肢体施加的每一自发影响，都使我们的土像仿佛通过某种机械装置活跃而可移动。因为最纯粹形式的热与最微妙形式的液体，通过混合与结合各自的力量而合一，以其水分滋养并维系大脑；进而，从该器官产生的蒸气被稀薄至最纯粹的状态，涂抹那包裹大脑的膜。此膜从上向下延伸，如管子般穿过连续的椎骨，它本身及其所含的骨髓与脊柱基部相接。它自身像御者一样，将冲动与力量赋予所有骨骼与关节的交接点以及肌肉的分支，以促成各部分的运动或静止。\par

25. 由于这个原因，在我看来，它也得到了更稳固的防御：在头部，它以双重的骨骼庇护环绕四周；在颈部的椎骨中，它凭借脊柱突起的壁垒以及椎骨形状本身的多样交织而受到保护，由此它免于一切伤害，借着环绕的防御享有安全。\par

26. 同样，人们可以认为心脏本身就像某座坚固的房屋，配备了最坚实的防御，由四周骨骼的围墙加固；因为后面有脊柱，两侧由肩胛骨加强，左右两侧肋骨的包围位置使它们之间的中部难以受伤；前面有胸骨和锁骨的结合作为防御，从而其安全在各方面都得到保护，免受外部危险之害。\par

27. 正如我们在农艺中所见，当雨从云中降下或河水泛滥使下方的土地饱含水分时（让我们假设一个园子，它在其范围之内滋养着无数种类的树木，以及从地里生长的一切植物形态，我们以极多细节的多样性沉思它们的形状、品质和特性）；那么，当这些植物在同一处被液体元素滋养时，供应给其中每一株的水分力量在性质上是一的；但被滋养的植物的特性将液体元素改变为不同的品质；因为同一物质在苦艾中变得苦涩，在毒参中变成致命的汁液，在其他植物中变得不同——在藏红花、香脂、罂粟中：在一处它变热，在另一处变冷，在另一处获得中间性质；在月桂和乳香中它散发香气，在无花果和梨中它变甜，通过葡萄藤它变成葡萄和酒；苹果的汁液、玫瑰的红色、百合的光辉、紫罗兰的蓝色、风信子染料的紫色，以及我们在地上所见的一切，都来自同一水分，并在形状、外表和品质上分为如此多的种类；同样的奇迹在我们存在的有生命土壤中由自然运作，或者更确切地说，由自然的主运作。骨、软骨、静脉、动脉、神经、韧带、肉、皮肤、脂肪、头发、腺、指甲、眼睛、鼻孔、耳朵——所有这类事物，以及无数其他事物，虽然由各种特性彼此区分，却由同一种滋养方式按它们本性的方式得到滋养，即当滋养物与任何实体密切接触时，它也随之按照它所接近的对象而改变，并适应并联合于该部分的特殊本性。因为如果它在眼睛附近，它就混合于视觉部分，并通过围绕眼球的不同膜层恰当地分配到各个部分；或者，如果它流向听觉部分，它就与听觉本性混合；若在唇中，它就变成唇；它在骨中变硬，在骨髓中变软，在肌腱中绷紧，在表面伸展，进入指甲，并经由相应的蒸发而精细化为头发的生长，产生多少有些卷曲或波状的头发，若它通过蜿蜒的通道行进；而若形成头发的蒸发途径是直的，它就使头发僵硬而直挺。\par

28. 然而，我们的论述已大大偏离其目的，深入自然之工，并试图描述我们的个别器官如何以及由什么材料形成——我指的是那些用于生命和良好生命的器官，以及我们最初划分中与这些一起包括的任何其他类别。\par

29. 因为我们的目的是要表明，我们构造的种子原因既不是无身体的灵魂，也不是无灵魂的身体，而是从有生命有身体的生物最初产生为有生命有灵魂的存在；我们的人性犹如乳母，用她自身拥有的资源接纳并培育它，它于是在两方面成长，并在相应部分显出其成长：——它立即通过这技巧性与科学性的形成过程，展示出交织于其中的灵魂力量，起初有些模糊，后来随着作品的完善而日益明亮。\par

30. 正如我们在石雕家身上所见——艺术家的目的是在石头上塑造某种动物的形象；他心中怀此目标，首先将石头同其同类的材料分离，然后通过凿去多余部分，从最初的轮廓逐步推进到心中所拟的模仿，即使外行的观察者也能从所见之物猜测其技艺的目标；再通过加工，使石头更接近心中对象的形貌；最后，在材料上造出完美而完成的形象，完成其技艺，而片刻之前还是无形之石的，变成了狮子、人，或艺术家所造的无论什么，不是通过将材料变为形象，而是通过将形象雕琢在材料上。若有人对灵魂作类似的设想，便离或然性不远了；因为我们说，那全智的自然，从同类的材料中取自人而来的部分，在自身内塑造雕像。正如随着石头的逐步加工，形象起初有些模糊，待作品完成后更为完美；同样，在其工具的塑造中，灵魂的形象表达在基质上，在不完美者中不完美，在完美者中完美；实在，若我们的本性未因邪恶而受损，它本应从起初就是完美的。因此，我们在那易受情感及具有动物性本性的受生中的共同之处，导致了天主的肖像并非在形成时立即发光，而是通过灵魂中属物质且更属受造动物的那些属性，以某种方法和顺序将人引至完美。\par

31. 伟大的宗徒在他的\WorkTitle{致格林多人书}中也教导我们某种类似的道理，他说：\BibleQuote{我作孩子的时候，说话像孩子，理解像孩子，思想像孩子；但当我成为成人时，我就把孩童的事丢弃了。}\BibleRef{格前 13:11}这并不是说，在成人中出现的灵魂与我们在孩童中所知的不同，孩童的理智消失而成人的理智在我们内生成；而是同一个灵魂，在孩童身上显示其不完善的状态，在成人身上显示其完善的状态。\par

32. 因为我们说，那些萌发生长的东西是有生命的，没有人会否认一切参与生命和自然运动的东西是富有生机的，但同时也不能说这样的生命分有完美的灵魂——因为，虽然某种生机的运作存在于植物中，它却不达到感觉的运动；另一方面，虽然某种更进一步的生命力存在于野兽中，它也不达到完美，因为它自身不含有理性与理智的恩宠。\par

33. 同样，我们说真正而完美的灵魂是人的灵魂，被每一种运作所认识；而其他分享生命的东西，我们因惯常的用语滥用而称之为有生机的，因为这些情况下灵魂并非以完美的状态存在，而只有灵魂运作的某些部分——在人类中（根据梅瑟关于人类起源的神秘记述）我们也得知，当人使自己类似这个感性的世界时，这些部分也附加上了。因此，保禄劝告那些能听他话的人抓住完美，也指示他们达到这一目的的方式，告诉他们说必须\Quote{脱去旧人}，穿上那\BibleQuote{按创造他者的形象更新的新人。}\BibleRef{哥 3:9}\par

34. 但愿我们都回到天主起初造人时所赐予的那神圣恩宠中，那时祂说：\Quote{让我们按我们的肖像和模样造人。}愿光荣与权能归于祂，至于无穷世。阿们。\par

\SourceRecord{Translated by H.A. Wilson.；Nicene and Post-Nicene Fathers, Second Series；Vol. 5.；Edited by Philip Schaff and Henry Wace.；Buffalo, NY: Christian Literature Publishing Co.,；1893.；<http://www.newadvent.org/fathers/2914.htm>.}
